%========================================
%. CONFIGURATION ARTICLE  COURS  SEC 101
%========================================
% 				CONTRIBUTIONS
%========================================

%************************************************************
% Chargement des variables du modèle
%************************************************************




%-------------------------------------------------------------
%               FR CYBERDEF SECOPS COURSE
%                          CONFIG NOTES File
%          SET BEAMER or Article/book Mode
%                    PRZMODE OR DOCMODE
%                             2020 eduf@ction
%-------------------------------------------------------------

% insert InterTitlle for each file compose main "Notes"
%\newcommand{\INTERTITLE}{OK}

 %------------------------------------------
 % Use global project CYBERDEF 
 %            Configuration file
 %------------------------------------------
%-----------------------------------------------
%               FR CYBERDEF SECOPS COURSE
%                         GLOBAL CONFIG  File
%                             2020 eduf@ction
%-----------------------------------------------

% just used for quick access to samples on project editing tool
%%-------------------------------------------------------------
%               FR CYBERDEF SECOPS COURSE
%              SAMPLES for edx & u  commands
%                             2020 eduf@ction
%-------------------------------------------------------------


% 												I M A G E S   
%-------------------------------------------------------------------------
%\upicture{../Latex/Sources/EDU/Pictures/img-liovar-incidents}{Incidents}{01}{lb:mtbf}
%-------------------------------------------------------------------------


%\mode<all>{\picframe{../Latex/Sources/EDU/Pictures/img-cycle}{Cycle de vie de gouvernance Cyberdef}{0.3}{lbl:cycle}}


% Planche PRZ 
% begin frame \mode<all>{\texframe{Déroulement}{présentation Gestion de la menace}
%{

%}} % end frame



%\begin{table}[!hbtp]
%\resizebox{\textwidth}{!}{%
%\begin{tabular}{ l|l l}
%%\begin{tabular}{l*{3}{S[table-format=2.1]}}
%\toprule
%\textbf{Package\_name }                 & \textbf{C\&C }                                        & \textbf{sha256}                                                            \\ \hline
%caracal.raceinspace.astronaut & http://api.lulquid.xyz                       & f1d32c17a169574369088...\\ \hline
%
%\bottomrule
%\end{tabular}
%}
%\caption{Exemple de  données de TI (Threat Intelligence)}
%\label{tab:ExempleTI}
%\end{table}

%----------------------------
% Path
\def\upath{../Latex}

%-------------------------------------------------------------
%			 FR CYBERDEF SECOPS COURSE
% 					INIT Commands File
% 		Select BEAMER or Article/book Mode
% 			with sub extension .doc or .prz
%                           2020 eduf@ction
%-------------------------------------------------------------
\RequirePackage{xstring}
\RequirePackage[realmainfile]{currfile}

\newcommand{\udocname}{\jobname}

\IfSubStr*{\currfilebase}{.prz}{\newcommand{\PRZMODE}{YES}}{\newcommand{\DOCMODE}{YES}}
   
\IfSubStr*{\currfilebase}{.edt}{\newcommand{\EDTMODE}{YES}}{\newcommand{\NOMODE}{YES}}

\input{../Latex/Templates/Magic/commons.inc/edx.library.tex}

\def\utemplatepath{../Latex/Templates/Magic}

\def\edxlibspath{\utemplatepath/edxlibs}
%\input{\edxlibspath/load.tex} à mettre dans les class article, book ...

\def\uauthorphoto{\utemplatepath/commons.inc/CommonsPictures/missionauthorphoto.png}


\def\utemplatetheme{cnam}
\def\ubeamermodel{edx}
\def\uproject{CNAMSEC101}


\newcommand{\upictureext}{pdf}

 %----------------------------
 % Basic Var
 
 %\setcounter{cntx}{\year+1}
\newcommand{\INFODistrib}{Notes de cours SECOPS 2022-2023}
\newcommand{\uJournalInfo}{Orange Cyber et Cloud avec CNAM Bretagne, Cybersécurité SEC101, edufaction}

\newcommand{\udescription}{La sécurité opérationnelle au coeur de la cyberdéfense d'entreprise}

\newcommand{\uUniversity}{Conservatoire National des Arts et Métiers}
\newcommand{\uCompany}{Orange Campus Cyber}
 %\newcommand{\uJournalInfo}{CNAM Bretagne, Cybersécurité SEC101, eduf@ction}
\newcommand{\uinstituteshort}{CNAM}
\newcommand{\uinstitute}{\uUniversity}
\newcommand{\uchaire}{Chaire de Cybersécurité}
\newcommand{\uversion}{Cnam Bretagne}





%----------------------------
% All other variables are defined in
 % file : *.lesson.tex
\newcommand{\ushorttitle}{CYBERDEF SEC101}
\newcommand{\uCoursetittle}{Eléments de Sécurité Opérationnelle}
\newcommand{\uMinilogo}{Cyberdéfense d'entreprise}
\newcommand{\uCourseLongName} {Eléments de sécurité opérationnelle en cyberdéfense d'entreprise}
%\newcommand{\utitle}{Cyberdéfense d'entreprise}

%----------------------------
% Book versus Course Notes
%----------------------------
% Use "e document' "e cours"
\newcommand{\edoc}{e document\xspace}
\newcommand{\ecours}{e cours\xspace} 
\newcommand{\etitle}{éléments de sécurité opérationnelle en cyberdéfense d'entreprise} 
\newcommand{\fichetech}{une fiche TECHNO}


%%----------------------------
%% Spécifique BOOK
% %----------------------------
\newcommand{\printer}{Publication limitée pour le \uUniversity \xspace et Orange Campus Cyber}

%%----------------------------
% Attention : Use hard coded Pictures Path
% (PicframeS)
\newcommand{\picpath}{../Latex/Sources/EDU/SRC1/Pictures}
%-------------------------------------------------------------

%%----------------------------
%% Spécifique COURSE
% %----------------------------

\newcommand {\parttitlecourse}{TITRE PARTIE (config)}
\newcommand{\coursenumber}{1}

\newcommand{\uGlossaryAcronyms}{
\loadglsentries{../Latex/Sources/commons/acronymes.tex}
\loadglsentries{../Latex/Sources/commons/glossaire.tex}
}

\newcommand{\ucourseresume}{\normalsize Aborder la sécurité des systèmes d’information sous l’angle d'une sécurité dynamique est un axe qui depuis quelques années apporte de nouvelle manière d’aborder la protection, la défense, et la résilience des systèmes d’information. La transformation digitale de l'entreprise modifie et rend plus flous les périmètres des systèmes d’informations. Cela nécessite une approche élargie du risque numérique et des nouvelles architectures de systèmes d'information intégrant des technologies de cloud et de cyberdéfense. Malgré la mise en place de mesures et de technologies de protection de plus en plus élaborées, l'impact d'une attaque ayant franchi ces frontières poreuses à considérablement augmenté.  Cette compilation des notes de cours élaborée dans le cadre d'un cours d'introduction à la cybersécurité et la cyberdéfense d'entreprise aborde les grands éléments fondamentaux permettant d'appréhender le domaine. Il permet d'en comprendre les enjeux, les codes, et acteurs. Protéger l'ensemble de l'entreprise alors qu'il est complexe de définir ses frontières est illusoire. Identifier les actifs essentiels ou vitaux et mettre en place les moyens adaptés à leur protection et leur défense est une démarche tactique qui permet de graduellement réduire ses cyber-risques tout en assurant un pilotage globale de la gouvernance de la sécurité sur ses volets de conformité.}

 %------------------------------------------
 
\renewcommand{\uMinilogo}{\includegraphics[width=0.04\paperwidth]{\utemplatepath/commons.inc/CommonsPictures/\utemplatetheme.theme/shield-20.pdf}}
 
\newcommand {\Ucontribute}
{
%%-------------------------------------------------------------
%               FR CYBERDEF SECOPS COURSE
%                            Contributions Rules
%                              2020 eduf@ction
%-------------------------------------------------------------

% \section{Contributions}

% \subsection{Comment contribuer}

Les notes et les présentations sont réalisées sous \hologo{LaTeX}. 

Vous pouvez contribuer au projet du cours CYBERDEF101. Les contributions peuvent se faire sous deux formes :

\begin{itemize}
  \item Corriger, amender, améliorer les notes publiées. A chaque session des modifications et évolutions sont apportées pour tenir compte des corrections de fond et de formes. 
  \item Ajouter, compléter, modifier des parties de notes sur la base de votre lecture du cours et de vos expertises dans chacun des domaines évoqués.
\end{itemize}

 Les fichiers sources sont publiés sur GITHUB dans l'espace : \ulink{https://github.com/edufaction/CYBERDEF101}{(edufaction/CYBERDEF101)}.
 
%%-------------------------------------------------------------
%               FR CYBERDEF SECOPS COURSE
%                Contributions DATAS & NAMES
%                              2020 eduf@ction
%-------------------------------------------------------------

\subsection{Les contributeurs/auteurs du cours}

\subsubsection{co-auteurs}

%-------------------------------------------------------------
%               FR CYBERDEF SECOPS COURSE
%                Contributions DATAS & NAMES
%                              2020 eduf@ction
%-------------------------------------------------------------
%           CO AUTEURS - Rédacteur de chapitre
%-------------------------------------------------------------

\coauthor{David  BATANY}{2019-2020}{Architecture et fonctionnement des Botnets}{Cnam SEC101}

\subsubsection{contributeurs}

%-------------------------------------------------------------
%               FR CYBERDEF SECOPS COURSE
%                Contributions DATAS & NAMES
%                              2020 eduf@ction
%-------------------------------------------------------------
%                 CONTRIBUTIONS AU COURS
%-------------------------------------------------------------


\coauthor{Céline JUBY}{2020}{Contributions d'amélioration et relectures}{Orange Cyberdefense}


} 

%------------------------------------------
 %     BEAMER OU ARTICLE MODEL
 %------------------------------------------
  
\ifdefined\PRZMODE 
	\providecommand{\umainload}{%-------------------------------------------------------------
%               FR CYBERDEF SECOPS COURSE
%                          BEAMER MAIN FILE
%                             2020 eduf@ction
%-------------------------------------------------------------

%\begin{frame}
%\frametitle<presentation>{TITRE DIAPO}
% frame text
%\end{frame}

%-------------------------------------------------------------
\documentclass[ignorenonframetext, 8pt, allowframebreaks]{beamer}

%-----------------------------------------
% Load Edxlibs in MAGIC path
%-----------------------------------------
% Need : \edxlibspath defined

\input{\edxlibspath/edxlib-placeholder.tex}
\input{\edxlibspath/edxlib-epigraphe.tex}
\input{\edxlibspath/edxlib-quizzmoodle.tex}


\usepackage[utf8]{inputenc}
\usepackage[T1]{fontenc}
\usepackage{csquotes}
\usepackage[french]{babel}
\usepackage{lmodern}
\usepackage{multicol}

\input{\utemplatepath/commons.inc/CommonsPictures/\utemplatetheme.theme/colors.inc} % set uMainColor

%-------------------------------------------------------------
%               FR CYBERDEF SECOPS COURSE
%                                  EDX PACKAGES
%.    Commons packages for BOOK, NOTES (.doc) 
%                            and BEAMER (.prz)
%                              2020 eduf@ction
%-------------------------------------------------------------

%\usepackage[hyphens,spaces,obeyspaces]{url}
%\usepackage[hyphens,spaces,obeyspaces]{url}



\PassOptionsToPackage{hyphens,spaces,obeyspaces}{url}

\usepackage[utf8]{inputenc}
\usepackage[french]{babel}
\usepackage[T1]{fontenc}
%\usepackage{markdown}



\usepackage{ccicons}

\usepackage{tikz} % Required for drawing custom shapes
\usetikzlibrary{tikzmark} 
\usetikzlibrary{mindmap,trees, backgrounds}


\usepackage{verbatim}
\usepackage{wrapfig}
\usepackage{xpatch}

\usepackage{hologo}

\usepackage[figure]{totalcount}

\usepackage{fontawesome}

\usepackage[skins]{tcolorbox} % don't forget SKINS option

\usepackage{listings}
\usepackage{upquote}

\lstset{
upquote=true,columns=flexible, basicstyle=\ttfamily,
language=HTML, 
frameround=tttt,
commentstyle=\color{gray},
identifierstyle=\color{blue},
keywordstyle=\color{uMainColor}\bfseries,
xleftmargin=2em,
xrightmargin=2em,
aboveskip=\topsep,
belowskip=\topsep, 
frame=single,
rulecolor=\color{uMainColor},
backgroundcolor=\color{uMainColor!5},
breaklines,
breakindent=1.5em,
showspaces=false,
showstringspaces=false,
showtabs=false,
xleftmargin=2em,
xrightmargin=2em,
aboveskip=1em,
belowskip=1em,
}



\usepackage{graphicx} % Required for including pictures
%\graphicspath{{Pictures/}} % Specifies the directory where pictures are stored

\usepackage{booktabs} % Required for nicer horizontal rules in tables

\usepackage{setspace}
\setstretch{1,2}

\usepackage{csquotes}

\usepackage{fancyhdr}

\usepackage{makeidx} 

\usepackage{xspace} 

\usepackage{floatflt} 

\usepackage{hyperref}
\usepackage{cleveref}

%pdfdisplaydoctitle={\utitle}, Attention le titre ne doit pas contenir de saut de ligne 

\hypersetup{
pdftitle={\uCoursetittle},
pdfauthor={\uauthor},
hidelinks,
backref=true,
pagebackref=true,
hyperindex=true,
pdfcenterwindow = true,
colorlinks=true,
breaklinks=true,
bookmarks=true,
allcolors=uMainColor,
pdfsubject = {\udescription},
pdfcreator = {MACTEX on TEXPAD},
bookmarksopen=true}


% urlcolor=cnam,
% runcolor=uMainColor,
%\pdfstringdefDisableCommands{%
%  \def\\{}%
%  \def\texttt#1{<#1>}%
%}

\usepackage[style=numeric,citestyle=numeric,sorting=nyt,sortcites=true,autopunct=true,hyperref=true,abbreviate=false,backref=true,backend=biber]{biblatex}

%\usepackage{cleveref} old position

\usepackage{multicol}

%--------------------------------------------------------
% pour la génération du QRCODE en page de garde d'article
\usepackage{qrcode} 

%--------------------------------------------------------
% POur la date d'impression en premiere page d'article
\usepackage[useregional]{datetime2}

%--------------------------------------------------------
% Gestion des glossaires et Sources/EDU/acronymes
% Un fichier acroymes.tex centralisé
\usepackage[acronym]{glossaries}

%\makeglossaries
\makenoidxglossaries 

\loadglsentries{Latex/Sources/commons/acronymes.tex}
\loadglsentries{Latex/Sources/commons/glossaire.tex}


%\renewcommand{\bibname}{R\'ef\'erences}

%--------------------------------------------------------
% Pour exclure des "environnement", ici techworkbox par exemple, travaux pour étudiants, non affiché dans le BOOK.
\usepackage{comment}


%--------------------------------------------------------
% pour filigramme sur les pages
\usepackage{eso-pic}

%--------------------------------------------------------
% for \gappto{\UrlBreaks}{\UrlOrds}
% rather than 
%\makeatletter
% \g@addto@macro{\UrlBreaks}{\UrlOrds}
% \makeatother

\usepackage{etoolbox}

\usepackage{breakurl}


\newcommand{\ucite}[1]{\cite{#1}}
%\newcommand{\ucite}[1]{REF BIB #1}
\newcommand{\ugls}[1]{\gls{#1}}
%\newcommand{\ugls}[1]{GLOSSAIRE (#1)}
%\usepackage[Rejne]{fncychap}

%\usepackage{pythontex} % need shell escape

%
%\makeatletter
%\DeclareRobustCommand{\afilename}[1]{%
% \begingroup
%  % \lstname seems to change hyphens into \textendash
%  \def\textendash{-}%
%  \afilename@parse{#1}%
%  \edef\afilename@base{\detokenize\expandafter{\afilename@base}}%
%  \texttt{\afilename@base.\afilename@ext}%
% \endgroup
%}
%\makeatother

%-------------------------------------------------------------
%               FR CYBERDEF SECOPS COURSE
%                                  EDX PACKAGES
%.    Commons COMMANDS for BOOK, NOTES (.doc) 
%                            and BEAMER (.prz)
%                              2020 eduf@ction
%-------------------------------------------------------------

\newcommand{\ucite}[1]{\cite{#1}}
%\newcommand{\ucite}[1]{REF BIB #1}
\newcommand{\ugls}[1]{\gls{#1}}
%\newcommand{\ugls}[1]{GLOSSAIRE (#1)}
\newcommand{\uac}[1]{\acrshort{#1}}



% \newcommand{\ExtractBeforeLastSlash}[1]{\directlua{extract_before_last_slash("#1")}}
% \newcommand{\ExtractAfterLastSlash}[1]{\directlua{extract_after_last_slash("#1")}}
% \newcommand{\RemoveExtension}[1]{\directlua{remove_extension("#1")}}
% \newcommand{\ExtractExtension}[1]{\directlua{extract_extension("#1")}}


\newcommand{\includepicturesubpath}[1]
{#1}

%{\ExtractBeforeLastSlash{#1}/\upicturethemepath\RemoveExtension{\ExtractAfterLastSlash{#1}}.\ExtractExtension{#1}}
    

%---------------------------------------------------
%										 	AUDIO NOTE 
%---------------------------------------------------

\newcommand\audionote[2]{%
  \note{#2}
  \immediate\write\audioexport{==AudioSlide=============}
  \immediate\write\audioexport{#1}
  \immediate\write\audioexport{\insertframenumber}
  \immediate\write\audioexport{---------------------------------}
  \immediate\write\audioexport{#2}
  \immediate\write\audioexport{---------------------------------}
}%

%---------------------------------------------------
%										 	TEXT STYLE SHORTCUT 
%---------------------------------------------------

\newcommand{\tcb}[1]{\textbf{\textcolor{uMainColor}{#1}}}

\newcommand{\voir}[1]{(Schéma~\ref{#1},~page~\pageref{#1}}

\newcommand{\graphicxextension}{pdf}

\newcommand\myemptypage{
   
    \thispagestyle{empty}
    \addtocounter{page}{-1}
    \newpage
  \null
    %CECI EST UNE PAGE VIDE
       \vfill
    }


\newcommand{\uNA}[2]
{
\newacronym{a#1}{#1}{#2}
}

\newcommand{\coauthor}[4]
{
\noindent (#2) \textbf{#1} - #4 : \textit{#3}
}


\newcommand\UKword[1]{\emph{#1}\xspace}
\newcommand\textUK[1]{\textit{\textbf{#1}}\xspace}
\newcommand\Ucenter[1]{\begin{center}{#1}\end{center}}

\newcommand\head[1]{\textbf{#1}}
\newcommand\tb[1]{\textbf{#1}}
\newcommand{\g}[1]{\og #1 \fg{}}
\newcommand{\uload}[1] {\input{../Latex/#1}}
\newcommand{\ucurcolor}{black}

\newcommand{\mak}[1]{\faicon{\aArrowCircleORight} #1}


%-----------------------------------------------------
%										 	\INFOGitHub 
%---------------------------------------------------


\newcommand{\GITfilename}{https://github.com/edufaction/CYBERDEF101/raw/main/Release/Builder/\jobname.pdf}

%\newcommand{\safeqrcode}[2][]{%
%\qrcode[#1]{\detokenize{#2}} }

\newcommand{\safeqrcode}[2][]{\qrcode[#1]{#2}} 

\newcommand{\INFOGitHub}
{
\begin{center} 
\setstretch{1.5}
 {Vérifiez la disponibilité d'une version plus récente de} \\
 {\ulink{\GITfilename}{\textbf{\jobname.pdf }sur GITHUB CYBERDEF} }   \\ 
{{\huge\ccbyncndeu}}  \\  
{Publication en \href{https://creativecommons.org/licenses/by-nc-nd/2.0/fr/}{Creative Common BY-NC-ND} by eduf@ction}    \\  % Copyright notice
{\safeqrcode[height=4 cm,padding]{\GITfilename}}  
\end{center} 
}

\newcommand{\INFOGitHubS}
{
\noindent{\safeqrcode[height=3 cm]{\GITfilename}} \\
\noindent {\tiny Téléchargez une version à jour} 
}



%-----------------------------------------------------
%										 	\upspicture 
%---------------------------------------------------
\newcommand{\upspicture}[3]
{
\renewcommand{\ucurcolor}{#3}
\resizebox{#2\textwidth}{!}{\input{../#1}}
}


%-----------------------------------------------------
%										 	\uindex 
%---------------------------------------------------


\newcommand{\uindex}[1]
{
#1\index{#1}\xspace
}

\newcommand{\ulindex}[1]{\index{#1}}

%-----------------------------------------------------
%										 	\includer 
%---------------------------------------------------

\newcommand{\includer}[1]
{
\input{../Latex/Sources/EDU/SRC1/Chapters/#1}
}

%-----------------------------------------------------
%										 	\edxdico 
%---------------------------------------------------

\newcommand{\edxdico}[2] {\textbf{#1}:\textit{#2} \index{#1}}

%-----------------------------------------------------
%  command DEF : InTextx
%-----------------------------------------------------
\newcommand\InTexta{no text}
\newcommand\InTextb{no text}
\newcommand\InTextc{no text}
\newcommand\rnc{\renewcommand}


%-----------------------------------------------------
%										 	ENV warningbox
%---------------------------------------------------

\newtcolorbox{warningbox}[2][]
{
  colframe = uMainColor!25,
  colback  = black!05,
  coltitle = uMainColor!20!black,
  fonttitle = \bfseries, 
  title    = #2,
  #1,
}


\newtcolorbox{techworkbox}[2][]{sharp corners,colframe = uMainColor!25,colback  = black!05,coltitle = uMainColor!20!black,fonttitle = \bfseries, title={\textcolor{uMainColor}{\faFile~\small {Sujet de mémoire SECOPS}}} : #2,#1,}


\newtcolorbox{notebox2}[2][]
{
sharp corners, 
colback = uMainColor!5!white, 
colframe = uMainColor!75!black,
fonttitle = \bfseries, 
colbacktitle= uMainColor!85!black,
title=#2,#1
}

\newtcolorbox{notebox}[2][]{colback=uMainColor!5!white,
colframe=uMainColor!75!black,fonttitle=\bfseries,
colbacktitle=uMainColor!85!black,enhanced,
attach boxed title to top right={yshift=+1mm},
title=#2,#1}


\newtcolorbox{toolsboxin}[2][]{sharp corners, colback=uMainColor!0!white,
colframe= black!25,fonttitle=\bfseries,
colbacktitle=black!60,enhanced,
attach boxed title to top left={yshift=+1mm},
title=#2,#1}


\newtcolorbox{techworkin}[2][]{sharp corners, colback=uMainColor!5!white,
colframe= uMainColor!100, fonttitle=\bfseries,
colbacktitle=black!5, enhanced,
 coltitle=black!50,
attach boxed title to top right ={yshift=+1mm},
title=#2,#1}



%-----------------------------------------------------
%										 	ENV ToolsBox
%---------------------------------------------------


\newcommand\tooleditor{No EDITOR in Datas.tex}
\newcommand\toolurl{No URL INSERTED in Datas.tex}
\newcommand\toolclass{No EDITOR in Datas.tex}
\newcommand\toolname{No NAME in Datas.tex}
\newcommand\toolanalyst{No ANALYST in Datas.tex}

\newcommand{\toolsbox}[2]
{
\input{../Latex/Sources/EDU/SRC1/SecComponents/#2.#1/datas.tex} 
\begin{toolsboxin}[sidebyside,righthand width=.7\textwidth]{#1}
\includegraphics[width=0.9\textwidth]{../Latex/Sources/EDU/SRC1/SecComponents/#2.#1/icon.png} \\
\tcblower
\input{../Latex/Sources/EDU/SRC1/SecComponents/#2.#1/description.tex}
\end{toolsboxin}
\noindent \small {\faCogs~Classe : \uppercase{\textbf{#2}}, Site de référence : \textbf{\ulink{\toolurl}{#1}}}
\newline
\small {\faGlobe~Editeur : \textbf{\tooleditor}}
\small {\faEye~Analyste : \textbf{\toolanalyst}}
}

%\renewcommand\toolsclass{SIEM}
%\renewcommand\toolsname{RSA Netwitness}
%\renewcommand\toolseditor{RSA}
%\renewcommand\toolsurl{https://www.rsa.com/fr-fr/products/threat-detection-response}

%-----------------------------------------------------
%										 	ENV  TechWork
%---------------------------------------------------


\newcommand{\techwork}[2]
{
\begin{techworkin} [righthand width=.8\textwidth]{\faCogs~\small {Fiche TECHNO} : #1}#2
\end{techworkin}
}

%-----------------------------------------------------
%										 	\uchap
%---------------------------------------------------
\ifdefined\INTERTITLE
\providecommand{\uchap}[1]
	{
	\begin{LARGE}
	\textbf{\Ucenter{#1 \\ -oOo-}}	
	\end{LARGE}
	}
\else
	\providecommand{\uchap}[1]{}
\fi

%-----------------------------------------------------
%											\uexpand
%---------------------------------------------------
\newcommand {\uexpand}[1]{
\mode<all>\input{../Latex/Sources/EDU/SRC1/Chapters/#1}
}


%-----------------------------------------------------
%											\spicture 
% direct include graphic with Scrivener
%-----------------------------------------------------

\newcommand{\spicture}[3]{
\begin{figure}[!hbtp] %h !hbtp
  \begin{center}
	 \includegraphics[width=0.8\textwidth]{#1}
  \end{center}
\caption{#3 - (slbl-#2)\label{slbl-#2}}
\end{figure}
}

%\StrBefore{#2}{.pdf}

%-----------------------------------------------------
%											\upicture 
%---------------------------------------------------

\ifdefined\PRZMODE

\providecommand{\upicture}[4]{
\framesubtitle{#2}
%\begin{figure}[!h] %h
  \begin{center}
	 \includegraphics[height=#3\textheight, width=#3\textwidth, keepaspectratio ]{\includepicturesubpath{#1.pdf}}
  \end{center}
%  \end{figure}
}

\else

\providecommand{\upicture}[4]{
\begin{figure}[!hbtp] %h !hbtp
  \begin{center}
	 \includegraphics[width=#3\textwidth]{\includepicturesubpath{#1.pdf}}
  \end{center}
\caption{#2 \label{#4}}
\end{figure}
}
	
\fi


\ifdefined\PRZMODE

\providecommand{\uwpicture}[4]{
\framesubtitle{#2}
%\begin{figure}[!h] %h
  \begin{center}
	 \includegraphics[height=0.8\textheight, width=0.8\textwidth, keepaspectratio ]{\includepicturesubpath{#1.pdf}}
  \end{center}
%  \end{figure}
}
\else
\providecommand{\uwpicture}[4]
{
\begin{wrapfigure}{O}{#3\textwidth} % inside / Outside
  \begin{center}
	 \includegraphics[width=#3\textwidth]{\includepicturesubpath{#1.pdf}}
\end{center}
\caption{#2 \label{#4}}
\end{wrapfigure}
}
\fi

%-----------------------------------------------------
%	\picframe for BEAMER (Image included in Beamer Slide)
%---------------------------------------------------

\ifdefined\PRZMODE

\providecommand{\picframe}[4]
{
\begin{frame}
\frametitle<presentation>{#2}
  \begin{center}

   \begin{tikzpicture}[remember picture, overlay]
    \node at (current page.center) {\includegraphics[width=#3\paperwidth, height=#3\paperheight,keepaspectratio]{\includepicturesubpath{#1.pdf}}};
    \end{tikzpicture}

  \end{center}
\end{frame}

}
\else
\providecommand{\picframe}[4]
{
\begin{figure}[!hbtp] % !hbtp ou !h 
\centering 
%{\includegraphics[width=#3\paperwidth, height=#3\paperheight,keepaspectratio]{\includepicturesubpath{#1.pdf}}}
\includegraphics[width=#3\textwidth]{\includepicturesubpath{#1.pdf}}
\caption{#2 \label{#4}}
\end{figure}
}
\fi

%-----------------------------------------------------
%.  for BEAMER (Image included in Beamer Slide)
%	\picframeS Simplified with Name and Label from One Param
%---------------------------------------------------


\ifdefined\PRZMODE

\providecommand{\picframeS}[2]
{
\begin{frame}
\frametitle<presentation>{#2}
  \begin{center}
	 \includegraphics[height=0.95\textheight, width=0.95\textwidth, keepaspectratio ]{\picpath/img-#1.\graphicxextension}
  \end{center}
\end{frame}

}
\else
\providecommand{\picframeS}[2]
{
\begin{figure}[!hbtp] % !hbtp ou !h 
\centering \includegraphics[width=0.9\textwidth]{\picpath/img-#1.\graphicxextension}
\caption{(s) #2 \label{lbl:#1}}
\end{figure}
}

\fi


%{figure}[!h]
%{wrapfigure}{O}{#3\textwidth}
%-----------------------------------------------------
%	FRAME			\wpicframe 
%---------------------------------------------------

\ifdefined\PRZMODE

\providecommand{\wpicframe}[4]
{
\begin{frame}
\frametitle<presentation>{#2}
  \begin{center}
	 \includegraphics[height=0.95\textheight, width=0.95\textwidth, keepaspectratio ]{#1.\graphicxextension}
  \end{center}
\end{frame}

}
\else
\providecommand{\wpicframe}[4]
{
\begin{wrapfigure}{O}{#3\textwidth}
\centering \includegraphics[width=#3\textwidth]{#1.\graphicxextension}
\caption{#2 \label{#4}}
\end{wrapfigure}
}
\fi

%-----------------------------------------------------
%	FRAME										\texframe 
%---------------------------------------------------
\ifdefined\PRZMODE

\providecommand{\texframe}[3]
{
\begin{frame}
\frametitle<presentation>{#1}
\framesubtitle<presentation>{#2}
#3
\end{frame}
}
\else
\providecommand{\texframe}[3]
{
#3
}

\fi

%-----------------------------------------------------
%											\uref 
%---------------------------------------------------
%\newcommand{\uref}[2]{(Voir~#1~\cref{#2} page~\pageref{#2})}

\newcommand{\uref}[2]{(Voir~#1~\ref{#2} page~\pageref{#2})}
%-----------------------------------------------------
%											\ubg
%---------------------------------------------------
\newcommand{\ubg}[1]{
\textbf{\g{\uppercase{#1}}}
}
%-----------------------------------------------------
%											\updfimage 
%---------------------------------------------------

\newcommand{\updfimage}[3]
{
\begin{wrapfigure}{R}{#3\textwidth}
  \begin{center}
	 \resizebox{#3\textwidth}{!}{\includegraphics{\includepicturesubpath{#1.pdf}}}
  \end{center}
\caption{#2}
\end{wrapfigure}
}


%************************PACKAGE***

%\usepackage{tikz}% Required for drawing custom shapes
\usetikzlibrary{shadows}


%-----------------------------------------------------
%											 		\rem 
% insertion d'une remarque avec R
%---------------------------------------------------
\newcommand\rem[1]{%
   \marginpar{%
   \tikzpicture[baseline={(title.base)}]
      \node[inner sep=5pt,text width=4cm,drop shadow={shadow yshift=-5pt,shadow xshift=5pt,uMainColor},fill=white] (box) {\vskip5pt \nointerlineskip #1};
      \node[right=10pt,inner sep=0pt,outer sep=10pt] (title) at (box.north west) {\bfseries\color{uMainColor}Remarque};
      \draw[draw=uMainColor,very thick](title.west)--(box.north west)--(box.south west)--(box.south east)--(box.north east)--(title.east);
      \fill[uMainColor]([yshift=-10pt]box.north west)--+(-5pt,-5pt)--+(0pt,-10pt);
   \endtikzpicture}%
}
%-----------------------------------------------------
%											 \utikzimage 
% insertion d'une image sur fichier tkz.tex
%---------------------------------------------------

\newcommand{\utikzimage}[3]
{
\begin{wrapfigure}{R}{#3\textwidth}
  \begin{center}
	 \resizebox{#3\textwidth}{!}{\input{../#1.tkz.tex}}
  \end{center}
\caption{#2}
\end{wrapfigure}
}

%-----------------------------------------------------
%														\link 
% Url et footnote du lien
%-----------------------------------------------------

% TT every character and hyphenate after it
\def\hyphenateAndTtWholeString #1{\xHyphenate#1$\wholeString\unskip}

\def\xHyphenate#1#2\wholeString {\if#1$%
    \else\transform{#1}%
    \takeTheRest#2\ofTheString\fi}

\def\takeTheRest#1\ofTheString\fi
{\fi \xHyphenate#1\wholeString}

\def\transform#1{\url{#1}\hskip 0pt plus 1pt}

% Define the \urlx command which works like \url, but with line brakes
\def\urlx #1{\href{#1}{\hyphenateAndTtWholeString{#1}}}
  

 %-----------------------------------------------------
%														\link 
%-----------------------------------------------------  
  
  \newcommand{\ulink}[2]
  {\href{#1}
  {#2~\raisebox{-0.2ex}{\faExternalLink}\footnote{\url{#1}}}\xspace}
  
  \newcommand{\wikipedia}[1]
  { \href{https://fr.wikipedia.org/wiki/#1}{\textcolor{grey}{#1}}\raisebox{0.5ex}{\tiny{\textcolor{uMainColor}{\faWikipediaW}}}}
  


 %-----------------------------------------------------
%														\utodo 
%-----------------------------------------------------

  \newcommand{\utodo}{\begin{center}{\noindent \textcolor{cyan}{\faGears~{En cours de rédaction, DRAFT non publiable} \faGears}}\end{center}}
  
    \newcommand{\utocomplete}{\begin{center}{\noindent \textcolor{blue}{\faClipboard~{à compléter, les éléments ne donnent qu'une vue trop réduite ou trop parcellaire du sujet}}}\end{center}}
  

 





%-------------------------------------------------------
\usepackage{\utemplatepath/LessonModel/Beamer/\ubeamermodel.model/edxbeamer}
%-------------------------------------------------------


% a charger après hyperref, pour avoir les guillemets sous beamer
\usepackage{ae}
\usepackage{aecompl}


\usepackage{environ}
\newwrite\audioexport



%\renewcommand{\familydefault}{pag}

%-------------------------------------------------------
% Optional: a subtitle to be displayed on the title slide

\title[\ushorttitle]{\lessonID~-~\utitle}
\subtitle{\tiny\uCourseLongName}

%-------------------------------------------------------
% The author(s) of the presentation:
%  - again first a short version to be displayed at the bottom;
%  - next the full list of authors, which may include contact information;

\author[\uauthorwriter]{ \LARGE{\uauthor} \\ \medskip {\small \url{\umaila} \and \url{\umailb}\\ {\small \url{\usitea}}}}

%-------------------------------------------------------
% The institute:
%  - to start the name of the university as displayed on the top of each slide
%    this can be adjusted such that you can also create a Dutch version
%  - next the institute information as displayed on the title slide

\institute[\uinstituteshort]{{\uinstitute \\ \uchaire}}


%-------------------------------------------------------
% Add a date and possibly the name of the event to the slides
%  - again first a short version to be shown at the bottom of each slide
%  - second the full date and event name for the title slide
\date[\uversion]{\tiny{Publication \INFODistrib ~du \\ \DTMnow}}


\setbeamertemplate{navigation symbols}{%
\insertslidenavigationsymbol
\insertframenavigationsymbol
\insertsubsectionnavigationsymbol
\insertsectionnavigationsymbol
\insertdocnavigationsymbol
\insertbackfindforwardnavigationsymbol
}

%\mode<all>{
%%-------------------------------------------------------------
%               FR CYBERDEF SECOPS COURSE
%
%                 DEFINITION DES Sources/EDU/acronymes

%                      $File :  Sources/EDU/acronymes.tex
%                             2020 eduf@ction
%-------------------------------------------------------------
% Require : glossaries package
%-------------------------------------------------------------
%\newacronym{XXX}{XXX}{XXX}
%-----------------------------------------------------------

%  Organismes

\newacronym{aCNIL}{CNIL}{Commission National Informatique et Liberté}
\newacronym{aOTAN}{OTAN}{Organisation du Traité de l'Atlantique Nord}
\newacronym{aANSSI}{ANSSI}{Agence Nationale de la Sécurité des systèmes d'information}
\newacronym{aCERT}{CERT}{Computer Emergency Response Team}
\newacronym{aAFNOR}{AFNOR}{Association Française de Normalisation}
\newacronym{aSIEM}{SIEM}{System Incident and Event Management}
\newacronym{aSIM}{SIM}{System Incident Management}
\newacronym{aSEM}{SEM}{System Event Management}
\newacronym{aIRP}{IRP}{Incident Response Plan, Plan de Réponse à Incident}
\newacronym{aISO}{ISO}{International Organization for Standardization}
\newacronym{aDICT}{DICT}{Disponibilité, Intégrité, Confidentialité, Traçabilité}
\newacronym{aPDCA}{PDCA}{Plan, Do, Check, Act (Roue de Deming)}
\newacronym{aGDPR}{GDPR}{General Data Protection Regulation ou RGPD Règlement général sur la protection des données}

%-----------------------------------------------------------

% ROLES ET METIERS

\newacronym{aOS}{OS}{Officier de sécurité}
\newacronym{aRSSI}{RSSI}{Responsable de la sécurité des systèmes d'information}
\newacronym{aDPO}{DPO}{Data Protection Officier}
\newacronym{aCIL}{CIL}{ Correspondant Informatique et Liberté}
\newacronym{aSSI}{SSI}{Sécurité des Systèmes d’information}
\newacronym{aGRC}{GRC}{Gouvernance Risques et Conformité}
\newacronym{aDTF}{DTF}{Defend The Flag}
\newacronym{aCTF}{CTF}{Capture The Flag}
\newacronym{aOIV}{OIV}{Opérateur d'Infrastructures Vitales}
\newacronym{aOSE}{OSE}{Opérateur de Services Essentiels}
\newacronym{aLPM}{LPM}{Loi de Programmation Militaire}
\newacronym{aEBIOS}{EBIOS}{Expression des Besoins et Identification des Objectifs de Sécurité}
\newacronym{aDLL}{DLL}{Dynamic Loads Library}
\newacronym{aACID}{ACID}{Authentification, Confidentialité, Intégrité, Disponibilité}

%-------------------------------------------------------
% Sources/EDU/acronymes DATABASE XLS Generation (cf. RESSOURCES DIR)
%-------------------------------------------------------
% DO NOT MODIFY DIRECTLY, USE XLS File (BEGIN)

\newacronym{a2FA}{2FA}{\textit{Two-Factor Authentication} : Authentification à deux facteurs}
\newacronym{aAAA}{AAA}{\textit{Authentication, Authorization, Accounting} : Authentification, autorisation, comptabilisation}
\newacronym{aAC}{AC}{\textit{Access Control} : Accès conditionnel}
\newacronym{aAES}{AES}{\textit{Advanced Encryption Standard} : Algorithme de chiffrement avancé}
\newacronym{aAPT}{APT}{\textit{Advanced Persistent Threat Group} : Groupe de cyber-espionnage avancé}
\newacronym{aASLR}{ASLR}{\textit{Address Space Layout Randomization} : Randomisation de l'espace d'adressage}
\newacronym{aAV}{AV}{\textit{Antivirus Software} : Logiciel antivirus}
\newacronym{aBOT}{BOT}{\textit{Robot (Malicious)} : Robot (malveillant)}
\newacronym{aCC}{C\&C}{\textit{Command And Control} : Commande et contrôle}
\newacronym{aCSRF}{CSRF}{\textit{Cross-Site Request Forgery Attack} : Attaque de confiance de site à site}
\newacronym{aDDOS}{DDOS}{\textit{Distributed Denial Of Service Attack} : Attaque de déni de service distribué}
\newacronym{aDES}{DES}{\textit{Data Encryption Standard} : Algorithme de chiffrement à clé de 64 bits}
\newacronym{aDLP}{DLP}{\textit{Data Loss Prevention} : Gestion de la protection des données sensibles}
\newacronym{aEDR}{EDR}{\textit{Endpoint Detection And Response} : Détection et réponse aux comportements anormaux d'un terminal}
\newacronym{aEK}{EK}{\textit{Exploit Kit} : Kit d'exploit}
\newacronym{aFTP}{FTP}{\textit{File Transfer Protocol} : Protocole de transfert de fichier}
\newacronym{aHTTPS}{HTTPS}{\textit{Protocol For Securing Communications On The Internet} : Protocole de sécurisation HTTP}
\newacronym{aIAM}{IAM}{\textit{Identity And Access Management} : Identité et Gestion des Accés}
\newacronym{aIDS}{IDS}{\textit{Intrusion Detection System} : Système de détection d'intrusion}
\newacronym{aIIS}{IIS}{\textit{Internet Information Services Server}}
\newacronym{aIPS}{IPS}{\textit{Intrusion Prevention System} : Système de prévention d'incident}
\newacronym{aISA}{ISA}{\textit{Internet Security And Acceleration Server} : Serveur de sécurité et d'accélération Internet (Microsoft)}
\newacronym{aLDAP}{LDAP}{\textit{Lightweight Directory Access Protocol} : Protocole de gestion d'annuaire léger}
\newacronym{aMD5}{MD5}{\textit{Message-Digest Algorithm 5}}
\newacronym{aMFA}{MFA}{\textit{Multi-Factor Authentication} : Authentification multi-facteurs}
\newacronym{aNAC}{NAC}{\textit{Network Access Control}}
\newacronym{aOAuth}{OAuth}{\textit{Open Authorization Protocol}}
\newacronym{aPAM}{PAM}{\textit{Privilege Access Management}}
\newacronym{aPKI}{PKI}{\textit{Public Key Infrastructure} : Infrastructure de clefs publiques}
\newacronym{aRSA}{RSA}{\textit{Rivest-Shamir-Adleman Encryption Algorithm}}
\newacronym{aSAML}{SAML}{\textit{Security Assertion Markup Language}}
\newacronym{aSMS}{SMS}{\textit{Short Message Service}}
\newacronym{aSQLI}{SQLI}{\textit{Sql Injection} : Injection de script de commande SQL}
\newacronym{aSSH}{SSH}{\textit{Secure Shell Protocol}}
\newacronym{aSSL}{SSL}{\textit{Protocol For Securing Communications On The Internet} : Protocole de sécurisation des communications sur Internet}
\newacronym{aTKIP}{TKIP}{\textit{Temporal Key Integrity Protocol}}
\newacronym{aTLS}{TLS}{\textit{Transport Layer Security}}
\newacronym{aURL}{URL}{\textit{Uniform Resource Locator}}
\newacronym{aVPN}{VPN}{\textit{Virtual Private Network} : Réseau privé virtuel}
\newacronym{aWAF}{WAF}{\textit{Web Application Firewall}}
\newacronym{aWPA}{WPA}{\textit{Wifi Protected Access}}
\newacronym{aXSS}{XSS}{\textit{Cross-Site Scripting Attack}}

% DO NOT MODIFY (END)
%-------------------------------------------------------




%%-------------------------------------------------------------
%          FR CYBERDEF SECOPS COURSE
%
%                        GLOSSAIRE
%
%                 $File :  Sources/EDU/glossaire.tex
%                    2020 eduf@ction
%-------------------------------------------------------------
% Require : glossaries package
%-------------------------------------------------------------

%\newglossaryentry{xxx}
%{
%name={xxx},
%description={xxx}
%}

\newglossaryentry{gLPM}
{
name={Loi de programmation militaire},
description={Une loi de programmation militaire (LPM), est une loi visant à établir une programmation pluriannuelle des dépenses et donc du budget que l'État français consacre à ses forces armées. Les lois de programmation militaire ont des durées d'application de quatre, cinq ou six ans. La cyberdéfense fait partie des axes importants des lois de programmation}
}

\newglossaryentry{gCTF}
{
name={Capture The Flag},
description={Challenge consistant à exploiter ses connaissance et les vulnérabilités affectant des logiciels  et systèmes numériques de manière à s’introduire sur des réseaux , des ordinateurs, pour récupérer (capturer) des preuves de l’intrusion réussie dénommées G\g{des drapeaux (\textit{Flags} en anglais)}.}
}

\newglossaryentry{gDTF}
{
name={Defend The Flag},
description={Challenge consistant à détecter, résister à des attaquants qui tentent de s'introduire ou sont introduits dans des systèmes numériques de manière à défendre le bien considéré comme le plus sensible convoité par l'attaquant dénommé G\g{drapeau (\textit{Flags} en anglais)}.}
}


\newglossaryentry{gROOTKIT}
{
name={Rootkit},
description={ou simplement « kit » (est un « outil de dissimulation d'activité » ou  « code furtif » qui est une sorte de trousse à outils pour un pirate informatique). Il est constitué  de techniques mises en œuvre par un ou plusieurs logiciels, dont le but est d'obtenir et de pérenniser un accès ceci avec en déployant des techniques de furtivité et de dissimulation.}
}

%-------------------------------------------------------
% GLS DATABASE XLS Generation (cf. RESSOURCES DIR)
%-------------------------------------------------------
% DO NOT MODIFY DIRECTLY, USE XLS File (BEGIN)

\newglossaryentry{gAffiliate}{name={\textit{Affiliates}},description={(Affiliés): Cybercriminels qui mènent des cyberattaques, parfois sophistiquées, par le biais de rançongiciels qui sont mis à disposition par des développeurs ou concepteurs. Les gains générés par les affiliés sont partagés avec les développeurs via un système de commissions.}}

% DO NOT MODIFY (END)
%-------------------------------------------------------


%}

%++++++++++++++++++++++++++++++++
\begin{document}



  %-------------------------------------------------------
  % AudioNote Open File
  %-------------------------------------------------------
\mode<all>{\immediate\openout\audioexport=\jobname.txt}

%================================
\begin{frame}[plain]
 					\begin{figure}
  					 \centering
  					 \begin{columns}[c]
  \begin{column}{0.4\textwidth}
  
    \begin{figure}
        \begin{flushleft}
  			\includegraphics[width=0.9\textwidth]{\utemplatepath/commons.inc/CommonsPictures/\utemplatetheme.theme/maincourselogo.pdf}
  \end{flushleft}
  \end{figure}
      \end{column}
  
  \begin{column}{0.6\textwidth}
    \begin{figure}
        \begin{flushright}
 		    \includegraphics[width=0.7\textwidth]{\utemplatepath/commons.inc/CommonsPictures/\utemplatetheme.theme/maininstitutelogo.pdf}
    \end{flushright}
  \end{figure}
  \end{column}
\end{columns}
   							  
   							  \end{figure}
							
				  				\titlepage
\end{frame}
%================================



%================================
\begin{frame}
	\frametitle{Abstract}
	\begin{techworkin} [righthand width=.8\textwidth]{\faCogs~\small {Hashtags} : \ukeywords}\uabstract
	\end{techworkin}
	\end{frame}
%================================

%================================
\begin{frame}
 \frametitle{Sommaire}
      \begin{multicols}{2}
         \tableofcontents[hideallsubsections]
     \end{multicols}
 \end{frame}
 %================================

%-------------------------------------------------------
%  BEAMER MAIN TEX configured with ubody
%-------------------------------------------------------
\mode<all>{\ubody}

%-------------------------------------------------------
% Last Frame
%-------------------------------------------------------
\begin{frame}
 \frametitle{des questions ?}
  \framesubtitle{contacter \umaila}
  					\begin{figure}
  					 \centering
   						 \includegraphics[width=0.2\textwidth]{\utemplatepath/commons.inc/CommonsPictures/\utemplatetheme.theme/mainlogo.pdf}\\
				  \end{figure}
				 	 \centering \textit{Tous les documents publiés dans le cadre de ce cours sont perfectibles, \\ne pas hésiter à m'envoyer vos remarques !}
 \end{frame}

%-------------------------------------------------------
% Contributions
%-------------------------------------------------------
\begin{frame}
 \frametitle{Contributions}
\begin{columns}
		\begin{column}{0.1\textwidth}

 					\begin{figure}
  						  \centering
   						 \includegraphics[width=0.9\textwidth]{../Latex/Sources/EDU/SRC1/Contribs/contribute/ContribPictures/github.pdf}\\
   						 	 \includegraphics[width=0.9\textwidth]{../Latex/Sources/EDU/SRC1/Contribs/contribute/ContribPictures/overleaf.pdf}
				  \end{figure}
		\end{column}
		\begin{column}{0.9\textwidth}
							%-------------------------------------------------------------
%               FR CYBERDEF SECOPS COURSE
%                            Contributions Rules
%                              2020 eduf@ction
%-------------------------------------------------------------

% \section{Contributions}

% \subsection{Comment contribuer}

Les notes et les présentations sont réalisées sous \hologo{LaTeX}. 

Vous pouvez contribuer au projet du cours CYBERDEF101. Les contributions peuvent se faire sous deux formes :

\begin{itemize}
  \item Corriger, amender, améliorer les notes publiées. A chaque session des modifications et évolutions sont apportées pour tenir compte des corrections de fond et de formes. 
  \item Ajouter, compléter, modifier des parties de notes sur la base de votre lecture du cours et de vos expertises dans chacun des domaines évoqués.
\end{itemize}

 Les fichiers sources sont publiés sur GITHUB dans l'espace : \ulink{https://github.com/edufaction/CYBERDEF101}{(edufaction/CYBERDEF101)}.
 
		\end{column}
\end{columns}
 \end{frame}
 
 %-------------------------------------------------------
% Link to download file and copyright
%-------------------------------------------------------
\begin{frame}
\frametitle{Mises à jour régulières}
\framesubtitle{Eduf@ction~\umaila}
%\INFOGitHub
{\begin{center} 
\setstretch{2}
 {Vérifiez la disponibilité d'une version plus récente de} \\
 {\ulink{\GITfilename}{\textbf{\udocname.pdf }sur GITHUB CYBERDEF} }   \\ 
{{\huge\ccbyncndeu}}  \\  
{\the\year~eduf@ction - Publication en Creative Common BY-NC-ND }    \\  % Copyright notice
{\safeqrcode[padding]{\GITfilename}}  
\end{center} }
\end{frame}


  %-------------------------------------------------------
  % AudioNote Open File
  %-------------------------------------------------------
  \mode<all>{\immediate\closeout\audioexport}


\end{document}




}
\else

\ifdefined\EDTMODE 
	\providecommand{\umainload}{%-------------------------------------------------------------
%               FR CYBERDEF SECOPS COURSE
%                          ARTICLE MAIN FILE
%                             2020 eduf@ction
%-------------------------------------------------------------

\documentclass[10pt,fleqn,twoside]{\utemplatepath/LessonModel/Article/edxstyle} 
%edxstyle
\usepackage{beamerarticle}

\usepackage{ae}
\usepackage[french]{babel}
\usepackage[utf8]{inputenc}
\usepackage[T1]{fontenc}
\usepackage{lmodern}




\usepackage{titletoc} % Required for manipulating the table of contents

\input{\utemplatepath/commons.inc/CommonsPictures/\utemplatetheme.theme/colors.inc}

%-------------------------------------------------------------
%               FR CYBERDEF SECOPS COURSE
%                                  EDX PACKAGES
%.    Commons packages for BOOK, NOTES (.doc) 
%                            and BEAMER (.prz)
%                              2020 eduf@ction
%-------------------------------------------------------------

%\usepackage[hyphens,spaces,obeyspaces]{url}
%\usepackage[hyphens,spaces,obeyspaces]{url}



\PassOptionsToPackage{hyphens,spaces,obeyspaces}{url}

\usepackage[utf8]{inputenc}
\usepackage[french]{babel}
\usepackage[T1]{fontenc}
%\usepackage{markdown}



\usepackage{ccicons}

\usepackage{tikz} % Required for drawing custom shapes
\usetikzlibrary{tikzmark} 
\usetikzlibrary{mindmap,trees, backgrounds}


\usepackage{verbatim}
\usepackage{wrapfig}
\usepackage{xpatch}

\usepackage{hologo}

\usepackage[figure]{totalcount}

\usepackage{fontawesome}

\usepackage[skins]{tcolorbox} % don't forget SKINS option

\usepackage{listings}
\usepackage{upquote}

\lstset{
upquote=true,columns=flexible, basicstyle=\ttfamily,
language=HTML, 
frameround=tttt,
commentstyle=\color{gray},
identifierstyle=\color{blue},
keywordstyle=\color{uMainColor}\bfseries,
xleftmargin=2em,
xrightmargin=2em,
aboveskip=\topsep,
belowskip=\topsep, 
frame=single,
rulecolor=\color{uMainColor},
backgroundcolor=\color{uMainColor!5},
breaklines,
breakindent=1.5em,
showspaces=false,
showstringspaces=false,
showtabs=false,
xleftmargin=2em,
xrightmargin=2em,
aboveskip=1em,
belowskip=1em,
}



\usepackage{graphicx} % Required for including pictures
%\graphicspath{{Pictures/}} % Specifies the directory where pictures are stored

\usepackage{booktabs} % Required for nicer horizontal rules in tables

\usepackage{setspace}
\setstretch{1,2}

\usepackage{csquotes}

\usepackage{fancyhdr}

\usepackage{makeidx} 

\usepackage{xspace} 

\usepackage{floatflt} 

\usepackage{hyperref}
\usepackage{cleveref}

%pdfdisplaydoctitle={\utitle}, Attention le titre ne doit pas contenir de saut de ligne 

\hypersetup{
pdftitle={\uCoursetittle},
pdfauthor={\uauthor},
hidelinks,
backref=true,
pagebackref=true,
hyperindex=true,
pdfcenterwindow = true,
colorlinks=true,
breaklinks=true,
bookmarks=true,
allcolors=uMainColor,
pdfsubject = {\udescription},
pdfcreator = {MACTEX on TEXPAD},
bookmarksopen=true}


% urlcolor=cnam,
% runcolor=uMainColor,
%\pdfstringdefDisableCommands{%
%  \def\\{}%
%  \def\texttt#1{<#1>}%
%}

\usepackage[style=numeric,citestyle=numeric,sorting=nyt,sortcites=true,autopunct=true,hyperref=true,abbreviate=false,backref=true,backend=biber]{biblatex}

%\usepackage{cleveref} old position

\usepackage{multicol}

%--------------------------------------------------------
% pour la génération du QRCODE en page de garde d'article
\usepackage{qrcode} 

%--------------------------------------------------------
% POur la date d'impression en premiere page d'article
\usepackage[useregional]{datetime2}

%--------------------------------------------------------
% Gestion des glossaires et Sources/EDU/acronymes
% Un fichier acroymes.tex centralisé
\usepackage[acronym]{glossaries}

%\makeglossaries
\makenoidxglossaries 

\loadglsentries{Latex/Sources/commons/acronymes.tex}
\loadglsentries{Latex/Sources/commons/glossaire.tex}


%\renewcommand{\bibname}{R\'ef\'erences}

%--------------------------------------------------------
% Pour exclure des "environnement", ici techworkbox par exemple, travaux pour étudiants, non affiché dans le BOOK.
\usepackage{comment}


%--------------------------------------------------------
% pour filigramme sur les pages
\usepackage{eso-pic}

%--------------------------------------------------------
% for \gappto{\UrlBreaks}{\UrlOrds}
% rather than 
%\makeatletter
% \g@addto@macro{\UrlBreaks}{\UrlOrds}
% \makeatother

\usepackage{etoolbox}

\usepackage{breakurl}


\newcommand{\ucite}[1]{\cite{#1}}
%\newcommand{\ucite}[1]{REF BIB #1}
\newcommand{\ugls}[1]{\gls{#1}}
%\newcommand{\ugls}[1]{GLOSSAIRE (#1)}
%\usepackage[Rejne]{fncychap}

%\usepackage{pythontex} % need shell escape

%
%\makeatletter
%\DeclareRobustCommand{\afilename}[1]{%
% \begingroup
%  % \lstname seems to change hyphens into \textendash
%  \def\textendash{-}%
%  \afilename@parse{#1}%
%  \edef\afilename@base{\detokenize\expandafter{\afilename@base}}%
%  \texttt{\afilename@base.\afilename@ext}%
% \endgroup
%}
%\makeatother

%-------------------------------------------------------------
%               FR CYBERDEF SECOPS COURSE
%                                  EDX PACKAGES
%.    Commons COMMANDS for BOOK, NOTES (.doc) 
%                            and BEAMER (.prz)
%                              2020 eduf@ction
%-------------------------------------------------------------

\newcommand{\ucite}[1]{\cite{#1}}
%\newcommand{\ucite}[1]{REF BIB #1}
\newcommand{\ugls}[1]{\gls{#1}}
%\newcommand{\ugls}[1]{GLOSSAIRE (#1)}
\newcommand{\uac}[1]{\acrshort{#1}}



% \newcommand{\ExtractBeforeLastSlash}[1]{\directlua{extract_before_last_slash("#1")}}
% \newcommand{\ExtractAfterLastSlash}[1]{\directlua{extract_after_last_slash("#1")}}
% \newcommand{\RemoveExtension}[1]{\directlua{remove_extension("#1")}}
% \newcommand{\ExtractExtension}[1]{\directlua{extract_extension("#1")}}


\newcommand{\includepicturesubpath}[1]
{#1}

%{\ExtractBeforeLastSlash{#1}/\upicturethemepath\RemoveExtension{\ExtractAfterLastSlash{#1}}.\ExtractExtension{#1}}
    

%---------------------------------------------------
%										 	AUDIO NOTE 
%---------------------------------------------------

\newcommand\audionote[2]{%
  \note{#2}
  \immediate\write\audioexport{==AudioSlide=============}
  \immediate\write\audioexport{#1}
  \immediate\write\audioexport{\insertframenumber}
  \immediate\write\audioexport{---------------------------------}
  \immediate\write\audioexport{#2}
  \immediate\write\audioexport{---------------------------------}
}%

%---------------------------------------------------
%										 	TEXT STYLE SHORTCUT 
%---------------------------------------------------

\newcommand{\tcb}[1]{\textbf{\textcolor{uMainColor}{#1}}}

\newcommand{\voir}[1]{(Schéma~\ref{#1},~page~\pageref{#1}}

\newcommand{\graphicxextension}{pdf}

\newcommand\myemptypage{
   
    \thispagestyle{empty}
    \addtocounter{page}{-1}
    \newpage
  \null
    %CECI EST UNE PAGE VIDE
       \vfill
    }


\newcommand{\uNA}[2]
{
\newacronym{a#1}{#1}{#2}
}

\newcommand{\coauthor}[4]
{
\noindent (#2) \textbf{#1} - #4 : \textit{#3}
}


\newcommand\UKword[1]{\emph{#1}\xspace}
\newcommand\textUK[1]{\textit{\textbf{#1}}\xspace}
\newcommand\Ucenter[1]{\begin{center}{#1}\end{center}}

\newcommand\head[1]{\textbf{#1}}
\newcommand\tb[1]{\textbf{#1}}
\newcommand{\g}[1]{\og #1 \fg{}}
\newcommand{\uload}[1] {\input{../Latex/#1}}
\newcommand{\ucurcolor}{black}

\newcommand{\mak}[1]{\faicon{\aArrowCircleORight} #1}


%-----------------------------------------------------
%										 	\INFOGitHub 
%---------------------------------------------------


\newcommand{\GITfilename}{https://github.com/edufaction/CYBERDEF101/raw/main/Release/Builder/\jobname.pdf}

%\newcommand{\safeqrcode}[2][]{%
%\qrcode[#1]{\detokenize{#2}} }

\newcommand{\safeqrcode}[2][]{\qrcode[#1]{#2}} 

\newcommand{\INFOGitHub}
{
\begin{center} 
\setstretch{1.5}
 {Vérifiez la disponibilité d'une version plus récente de} \\
 {\ulink{\GITfilename}{\textbf{\jobname.pdf }sur GITHUB CYBERDEF} }   \\ 
{{\huge\ccbyncndeu}}  \\  
{Publication en \href{https://creativecommons.org/licenses/by-nc-nd/2.0/fr/}{Creative Common BY-NC-ND} by eduf@ction}    \\  % Copyright notice
{\safeqrcode[height=4 cm,padding]{\GITfilename}}  
\end{center} 
}

\newcommand{\INFOGitHubS}
{
\noindent{\safeqrcode[height=3 cm]{\GITfilename}} \\
\noindent {\tiny Téléchargez une version à jour} 
}



%-----------------------------------------------------
%										 	\upspicture 
%---------------------------------------------------
\newcommand{\upspicture}[3]
{
\renewcommand{\ucurcolor}{#3}
\resizebox{#2\textwidth}{!}{\input{../#1}}
}


%-----------------------------------------------------
%										 	\uindex 
%---------------------------------------------------


\newcommand{\uindex}[1]
{
#1\index{#1}\xspace
}

\newcommand{\ulindex}[1]{\index{#1}}

%-----------------------------------------------------
%										 	\includer 
%---------------------------------------------------

\newcommand{\includer}[1]
{
\input{../Latex/Sources/EDU/SRC1/Chapters/#1}
}

%-----------------------------------------------------
%										 	\edxdico 
%---------------------------------------------------

\newcommand{\edxdico}[2] {\textbf{#1}:\textit{#2} \index{#1}}

%-----------------------------------------------------
%  command DEF : InTextx
%-----------------------------------------------------
\newcommand\InTexta{no text}
\newcommand\InTextb{no text}
\newcommand\InTextc{no text}
\newcommand\rnc{\renewcommand}


%-----------------------------------------------------
%										 	ENV warningbox
%---------------------------------------------------

\newtcolorbox{warningbox}[2][]
{
  colframe = uMainColor!25,
  colback  = black!05,
  coltitle = uMainColor!20!black,
  fonttitle = \bfseries, 
  title    = #2,
  #1,
}


\newtcolorbox{techworkbox}[2][]{sharp corners,colframe = uMainColor!25,colback  = black!05,coltitle = uMainColor!20!black,fonttitle = \bfseries, title={\textcolor{uMainColor}{\faFile~\small {Sujet de mémoire SECOPS}}} : #2,#1,}


\newtcolorbox{notebox2}[2][]
{
sharp corners, 
colback = uMainColor!5!white, 
colframe = uMainColor!75!black,
fonttitle = \bfseries, 
colbacktitle= uMainColor!85!black,
title=#2,#1
}

\newtcolorbox{notebox}[2][]{colback=uMainColor!5!white,
colframe=uMainColor!75!black,fonttitle=\bfseries,
colbacktitle=uMainColor!85!black,enhanced,
attach boxed title to top right={yshift=+1mm},
title=#2,#1}


\newtcolorbox{toolsboxin}[2][]{sharp corners, colback=uMainColor!0!white,
colframe= black!25,fonttitle=\bfseries,
colbacktitle=black!60,enhanced,
attach boxed title to top left={yshift=+1mm},
title=#2,#1}


\newtcolorbox{techworkin}[2][]{sharp corners, colback=uMainColor!5!white,
colframe= uMainColor!100, fonttitle=\bfseries,
colbacktitle=black!5, enhanced,
 coltitle=black!50,
attach boxed title to top right ={yshift=+1mm},
title=#2,#1}



%-----------------------------------------------------
%										 	ENV ToolsBox
%---------------------------------------------------


\newcommand\tooleditor{No EDITOR in Datas.tex}
\newcommand\toolurl{No URL INSERTED in Datas.tex}
\newcommand\toolclass{No EDITOR in Datas.tex}
\newcommand\toolname{No NAME in Datas.tex}
\newcommand\toolanalyst{No ANALYST in Datas.tex}

\newcommand{\toolsbox}[2]
{
\input{../Latex/Sources/EDU/SRC1/SecComponents/#2.#1/datas.tex} 
\begin{toolsboxin}[sidebyside,righthand width=.7\textwidth]{#1}
\includegraphics[width=0.9\textwidth]{../Latex/Sources/EDU/SRC1/SecComponents/#2.#1/icon.png} \\
\tcblower
\input{../Latex/Sources/EDU/SRC1/SecComponents/#2.#1/description.tex}
\end{toolsboxin}
\noindent \small {\faCogs~Classe : \uppercase{\textbf{#2}}, Site de référence : \textbf{\ulink{\toolurl}{#1}}}
\newline
\small {\faGlobe~Editeur : \textbf{\tooleditor}}
\small {\faEye~Analyste : \textbf{\toolanalyst}}
}

%\renewcommand\toolsclass{SIEM}
%\renewcommand\toolsname{RSA Netwitness}
%\renewcommand\toolseditor{RSA}
%\renewcommand\toolsurl{https://www.rsa.com/fr-fr/products/threat-detection-response}

%-----------------------------------------------------
%										 	ENV  TechWork
%---------------------------------------------------


\newcommand{\techwork}[2]
{
\begin{techworkin} [righthand width=.8\textwidth]{\faCogs~\small {Fiche TECHNO} : #1}#2
\end{techworkin}
}

%-----------------------------------------------------
%										 	\uchap
%---------------------------------------------------
\ifdefined\INTERTITLE
\providecommand{\uchap}[1]
	{
	\begin{LARGE}
	\textbf{\Ucenter{#1 \\ -oOo-}}	
	\end{LARGE}
	}
\else
	\providecommand{\uchap}[1]{}
\fi

%-----------------------------------------------------
%											\uexpand
%---------------------------------------------------
\newcommand {\uexpand}[1]{
\mode<all>\input{../Latex/Sources/EDU/SRC1/Chapters/#1}
}


%-----------------------------------------------------
%											\spicture 
% direct include graphic with Scrivener
%-----------------------------------------------------

\newcommand{\spicture}[3]{
\begin{figure}[!hbtp] %h !hbtp
  \begin{center}
	 \includegraphics[width=0.8\textwidth]{#1}
  \end{center}
\caption{#3 - (slbl-#2)\label{slbl-#2}}
\end{figure}
}

%\StrBefore{#2}{.pdf}

%-----------------------------------------------------
%											\upicture 
%---------------------------------------------------

\ifdefined\PRZMODE

\providecommand{\upicture}[4]{
\framesubtitle{#2}
%\begin{figure}[!h] %h
  \begin{center}
	 \includegraphics[height=#3\textheight, width=#3\textwidth, keepaspectratio ]{\includepicturesubpath{#1.pdf}}
  \end{center}
%  \end{figure}
}

\else

\providecommand{\upicture}[4]{
\begin{figure}[!hbtp] %h !hbtp
  \begin{center}
	 \includegraphics[width=#3\textwidth]{\includepicturesubpath{#1.pdf}}
  \end{center}
\caption{#2 \label{#4}}
\end{figure}
}
	
\fi


\ifdefined\PRZMODE

\providecommand{\uwpicture}[4]{
\framesubtitle{#2}
%\begin{figure}[!h] %h
  \begin{center}
	 \includegraphics[height=0.8\textheight, width=0.8\textwidth, keepaspectratio ]{\includepicturesubpath{#1.pdf}}
  \end{center}
%  \end{figure}
}
\else
\providecommand{\uwpicture}[4]
{
\begin{wrapfigure}{O}{#3\textwidth} % inside / Outside
  \begin{center}
	 \includegraphics[width=#3\textwidth]{\includepicturesubpath{#1.pdf}}
\end{center}
\caption{#2 \label{#4}}
\end{wrapfigure}
}
\fi

%-----------------------------------------------------
%	\picframe for BEAMER (Image included in Beamer Slide)
%---------------------------------------------------

\ifdefined\PRZMODE

\providecommand{\picframe}[4]
{
\begin{frame}
\frametitle<presentation>{#2}
  \begin{center}

   \begin{tikzpicture}[remember picture, overlay]
    \node at (current page.center) {\includegraphics[width=#3\paperwidth, height=#3\paperheight,keepaspectratio]{\includepicturesubpath{#1.pdf}}};
    \end{tikzpicture}

  \end{center}
\end{frame}

}
\else
\providecommand{\picframe}[4]
{
\begin{figure}[!hbtp] % !hbtp ou !h 
\centering 
%{\includegraphics[width=#3\paperwidth, height=#3\paperheight,keepaspectratio]{\includepicturesubpath{#1.pdf}}}
\includegraphics[width=#3\textwidth]{\includepicturesubpath{#1.pdf}}
\caption{#2 \label{#4}}
\end{figure}
}
\fi

%-----------------------------------------------------
%.  for BEAMER (Image included in Beamer Slide)
%	\picframeS Simplified with Name and Label from One Param
%---------------------------------------------------


\ifdefined\PRZMODE

\providecommand{\picframeS}[2]
{
\begin{frame}
\frametitle<presentation>{#2}
  \begin{center}
	 \includegraphics[height=0.95\textheight, width=0.95\textwidth, keepaspectratio ]{\picpath/img-#1.\graphicxextension}
  \end{center}
\end{frame}

}
\else
\providecommand{\picframeS}[2]
{
\begin{figure}[!hbtp] % !hbtp ou !h 
\centering \includegraphics[width=0.9\textwidth]{\picpath/img-#1.\graphicxextension}
\caption{(s) #2 \label{lbl:#1}}
\end{figure}
}

\fi


%{figure}[!h]
%{wrapfigure}{O}{#3\textwidth}
%-----------------------------------------------------
%	FRAME			\wpicframe 
%---------------------------------------------------

\ifdefined\PRZMODE

\providecommand{\wpicframe}[4]
{
\begin{frame}
\frametitle<presentation>{#2}
  \begin{center}
	 \includegraphics[height=0.95\textheight, width=0.95\textwidth, keepaspectratio ]{#1.\graphicxextension}
  \end{center}
\end{frame}

}
\else
\providecommand{\wpicframe}[4]
{
\begin{wrapfigure}{O}{#3\textwidth}
\centering \includegraphics[width=#3\textwidth]{#1.\graphicxextension}
\caption{#2 \label{#4}}
\end{wrapfigure}
}
\fi

%-----------------------------------------------------
%	FRAME										\texframe 
%---------------------------------------------------
\ifdefined\PRZMODE

\providecommand{\texframe}[3]
{
\begin{frame}
\frametitle<presentation>{#1}
\framesubtitle<presentation>{#2}
#3
\end{frame}
}
\else
\providecommand{\texframe}[3]
{
#3
}

\fi

%-----------------------------------------------------
%											\uref 
%---------------------------------------------------
%\newcommand{\uref}[2]{(Voir~#1~\cref{#2} page~\pageref{#2})}

\newcommand{\uref}[2]{(Voir~#1~\ref{#2} page~\pageref{#2})}
%-----------------------------------------------------
%											\ubg
%---------------------------------------------------
\newcommand{\ubg}[1]{
\textbf{\g{\uppercase{#1}}}
}
%-----------------------------------------------------
%											\updfimage 
%---------------------------------------------------

\newcommand{\updfimage}[3]
{
\begin{wrapfigure}{R}{#3\textwidth}
  \begin{center}
	 \resizebox{#3\textwidth}{!}{\includegraphics{\includepicturesubpath{#1.pdf}}}
  \end{center}
\caption{#2}
\end{wrapfigure}
}


%************************PACKAGE***

%\usepackage{tikz}% Required for drawing custom shapes
\usetikzlibrary{shadows}


%-----------------------------------------------------
%											 		\rem 
% insertion d'une remarque avec R
%---------------------------------------------------
\newcommand\rem[1]{%
   \marginpar{%
   \tikzpicture[baseline={(title.base)}]
      \node[inner sep=5pt,text width=4cm,drop shadow={shadow yshift=-5pt,shadow xshift=5pt,uMainColor},fill=white] (box) {\vskip5pt \nointerlineskip #1};
      \node[right=10pt,inner sep=0pt,outer sep=10pt] (title) at (box.north west) {\bfseries\color{uMainColor}Remarque};
      \draw[draw=uMainColor,very thick](title.west)--(box.north west)--(box.south west)--(box.south east)--(box.north east)--(title.east);
      \fill[uMainColor]([yshift=-10pt]box.north west)--+(-5pt,-5pt)--+(0pt,-10pt);
   \endtikzpicture}%
}
%-----------------------------------------------------
%											 \utikzimage 
% insertion d'une image sur fichier tkz.tex
%---------------------------------------------------

\newcommand{\utikzimage}[3]
{
\begin{wrapfigure}{R}{#3\textwidth}
  \begin{center}
	 \resizebox{#3\textwidth}{!}{\input{../#1.tkz.tex}}
  \end{center}
\caption{#2}
\end{wrapfigure}
}

%-----------------------------------------------------
%														\link 
% Url et footnote du lien
%-----------------------------------------------------

% TT every character and hyphenate after it
\def\hyphenateAndTtWholeString #1{\xHyphenate#1$\wholeString\unskip}

\def\xHyphenate#1#2\wholeString {\if#1$%
    \else\transform{#1}%
    \takeTheRest#2\ofTheString\fi}

\def\takeTheRest#1\ofTheString\fi
{\fi \xHyphenate#1\wholeString}

\def\transform#1{\url{#1}\hskip 0pt plus 1pt}

% Define the \urlx command which works like \url, but with line brakes
\def\urlx #1{\href{#1}{\hyphenateAndTtWholeString{#1}}}
  

 %-----------------------------------------------------
%														\link 
%-----------------------------------------------------  
  
  \newcommand{\ulink}[2]
  {\href{#1}
  {#2~\raisebox{-0.2ex}{\faExternalLink}\footnote{\url{#1}}}\xspace}
  
  \newcommand{\wikipedia}[1]
  { \href{https://fr.wikipedia.org/wiki/#1}{\textcolor{grey}{#1}}\raisebox{0.5ex}{\tiny{\textcolor{uMainColor}{\faWikipediaW}}}}
  


 %-----------------------------------------------------
%														\utodo 
%-----------------------------------------------------

  \newcommand{\utodo}{\begin{center}{\noindent \textcolor{cyan}{\faGears~{En cours de rédaction, DRAFT non publiable} \faGears}}\end{center}}
  
    \newcommand{\utocomplete}{\begin{center}{\noindent \textcolor{blue}{\faClipboard~{à compléter, les éléments ne donnent qu'une vue trop réduite ou trop parcellaire du sujet}}}\end{center}}
  

 




%-------------------------------------------------------------
%               FR CYBERDEF SECOPS COURSE
%
%                                    EDX STYLE
%
%                              2020 eduf@ction
%-------------------------------------------------------------
\usepackage{enumitem}

\usepackage{parskip}
\setlength{\parindent}{0em}
\setlength{\parskip}{5pt} % 1ex plus 0.5ex minus 0.2ex}

\def\uCoverColor{}

%---------------------------------------------------------------
%	MARGINS
%---------------------------------------------------------------

\usepackage{geometry}

\geometry{
	paper=a4paper, 
%	top=1cm, 
%	bottom=1cm, 
%	left=1cm, 
%	right=1cm, 
%	headheight=20pt, 
%	footskip=1.4cm, 
%	headsep=10pt, 
%	showframe, 
	showcrop 
}

%---------------------------------------------------------------
%	FOOTNOTE (format)
%---------------------------------------------------------------

\makeatletter
\long\def\@makefntextFB#1{%
    \ifx\thefootnote\ftnISsymbol
        \@makefntextORI{#1}%
    \else
        \rule\z@\footnotesep
        \setbox\@tempboxa\hbox{\@thefnmark}%
            \ifdim\wd\@tempboxa>\z@
                \kern2em\llap{\@thefnmark.\kern0.5em}%
            \fi
        \hangindent2em\hangafter\@ne#1
    \fi}
\makeatother

%-----------------------------------------------------
%	COLOR
%-----------------------------------------------------
%\usepackage{xcolor} % Required for specifying colors by name
%\definecolor{uMainColor}{RGB}{160,0,0} 
\colorlet{edxcolorcover}{uCoverColor}
%\definecolor{edxcolorcover}{RGB}{160,0,0} 

%\definecolor{cnam}{RGB}{128, 0, 32} 
                 

  \definecolor{comments}{rgb}{0.7,0,0}    % rouge foncé
  \definecolor{link}{rgb}{0,0.4,0.6}      % ~RoyalBlue de dvips
  \definecolor{url}{rgb}{0.6,0,0}         % rouge-brun
  \definecolor{citation}{rgb}{0,0.5,0}    % vert foncé
  \definecolor{ULlinkcolor}{rgb}{0,0,0.3} % de ulthese.cls
  \definecolor{rouge}{rgb}{0.85,0,0.07}   % rouge bandeau identitaire
  \definecolor{or}{rgb}{1,0.8,0}          % or bandeau identitaire

%-----------------------------------------------------
%	FONTS
%-----------------------------------------------------

%%\usepackage{avant} % Use the Avantgarde font for headings
%\usepackage{times} % Use the Times font for headings
%\usepackage{mathptmx} % Use the Adobe Times Rgman as the default text font together with math symbols from the Sym­bol, Chancery and Com­puter Modern fonts

%%\usepackage{microtype} % Slightly tweak font spacing for aesthetics
%%\usepackage[utf8]{inputenc} % Required for including letters with accents
%\usepackage[T1]{fontenc} % Use 8-bit encoding that has 256 glyphs

% Using AVANT GARDE Family

%\renewcommand{\familydefault}{\sfdefault}

%\fontfamily{pag}\selectfont
%\renewcommand{\familydefault}{pag} % sfdefault pag lmss

%bch         Charter
%lmr         Latin Modern Roman
%lmss        Latin Modern Sans Serif
%lmssq       Latin Modern Sans Serif extended
%lmtt        Latin Modern Typewriter
%lmvtt       Latin Modern Typewriter proportional
%pag         Avant Garde
%pbk         Bookman
%pcr         Courier
%phv         Helvetica
%pnc         New Century Schoolbook
%ppl         Palatino
%ptm         Times
%put         Utopia
 
%\renewcommand{\rmdefault}{pag} % text 
%\renewcommand{\sfdefault}{pag} % titre


%\renewcommand{\rmdefault}{pag} % text 
%\renewcommand{\sfdefault}{phv} % titre


%-----------------------------------------------------
%	BIBLIOGRAPHY AND INDEX
%-----------------------------------------------------


\usepackage{calc} % For simpler calculation - used for spacing the index letter headings correctly

\usepackage{imakeidx}
%\makeindex % Tells LaTeX to create the files required for indexing

\makeindex[columns=2,intoc=true,options={-s \upath/indexstyle.ist}]

%-----------------------------------------------------
%	HEADERS AND FOOTERS
%-----------------------------------------------------


\usepackage{fancyhdr} % Required for header and footer configuration
\pagestyle{fancy} % Enable the custom headers and footers

\ifcsname chapter \endcsname   % verify if command name exists in the CLASS (Book, report, article ...)
% BEGIN IF
\renewcommand{\chaptermark}[1]{\markboth{\sffamily\normalsize\bfseries\chaptername\ \thechapter.\ #1}{}} % Styling for the current chapter in the header
  \else
%nothing
\fi
% END IF

\renewcommand{\sectionmark}[1]{\markright{\sffamily\normalsize\thesection\hspace{5pt}#1}{}} % Styling for the current section in the header

\fancyhf{} % Clear default headers and footers
\fancyhead[LE,RO]{\sffamily\normalsize\thepage} % Styling for the page number in the header
\fancyhead[LO]{\sffamily\normalsize\rightmark} % Print the nearest section name on the left side of odd pages
\fancyhead[RE]{\sffamily\normalsize\leftmark} % Print the current chapter name on the right side of even pages
\fancyfoot[RE,LO]{\sffamily\normalsize\uCoursetittle} % Uncomment to include a footer
\fancyfoot[LE,RO]{\hyperlink{toc}{\sffamily\normalsize\uMinilogo}} % Uncomment to include a footer
% \hyperlink{toc}
\renewcommand{\headrulewidth}{1.0pt} % Thickness of the rule under the header

\fancypagestyle{plain}{% Style for when a plain pagestyle is specified
	\fancyhead{}\renewcommand{\headrulewidth}{0pt}} 

% Removes the header from odd empty pages at the end of chapters
\makeatletter
\renewcommand{\cleardoublepage}{
\clearpage\ifodd\c@page\else
\hbox{}
\vspace*{\fill}
\thispagestyle{empty}
\newpage
\fi}


%----------------------------------------------------------------------------------------
%	MAIN TABLE OF CONTENTS
%----------------------------------------------------------------------------------------
%\usepackage{titletoc} % Required for manipulating the table of contents

\contentsmargin{0cm} % Removes the default margin

% Part text styling (this is mostly taken care of in the PART HEADINGS section of this file)
\titlecontents{part}
	[0cm] % Left indentation
	{\addvspace{20pt}\bfseries} % Spacing and font options for parts
	{}
	{}
	{}

% Chapter text styling
\titlecontents{chapter}
	[1.25cm] % Left indentation
	{\addvspace{12pt}\large\sffamily\bfseries} % Spacing and font options for chapters
	{\color{uMainColor!60}\contentslabel[\Large\thecontentslabel]{1.25cm}\color{uMainColor}} % Formatting of numbered sections of this type
	{\color{uMainColor}} % Formatting of numberless sections of this type
	{\color{uMainColor!60}\normalsize\;\titlerule*[.5pc]{.}\;\thecontentspage} % Formatting of the filler to the right of the heading and the page number

% Section text styling
\titlecontents{section}
	[1.25cm] % Left indentation
	{\addvspace{3pt}\sffamily\bfseries} % Spacing and font options for sections
	{\contentslabel[\thecontentslabel]{1.25cm}} % Formatting of numbered sections of this type
	{} % Formatting of numberless sections of this type
	{\hfill\color{black}\thecontentspage} % Formatting of the filler to the right of the heading and the page number

% Subsection text styling
\titlecontents{subsection}
	[1.25cm] % Left indentation
	{\addvspace{1pt}\sffamily\small} % Spacing and font options for subsections
	{\contentslabel[\thecontentslabel]{1.25cm}} % Formatting of numbered sections of this type
	{} % Formatting of numberless sections of this type
	{\ \titlerule*[.5pc]{.}\;\thecontentspage} % Formatting of the filler to the right of the heading and the page number

%\titlecontents{subsubsection}
%	[1.25cm] % Left indentation
%	{\addvspace{1pt}\sffamily\tiny} % Spacing and font options for subsections
%	{\contentslabel[\thecontentslabel]{1.25cm}} % Formatting of numbered sections of this type
%	{} % Formatting of numberless sections of this type
%	{\ \titlerule*[.5pc]{.}\;\thecontentspage} % Formatting of the filler to the right of the heading and the page number

\titlecontents{subsubsection}
[1.25cm]
  {\addvspace{1pt}\sffamily\tiny}
 {\contentslabel[\thecontentslabel]{1.25cm}}
  {}
  {}

% Figure text styling
\titlecontents{figure}
	[1.25cm] % Left indentation
	{\addvspace{1pt}\sffamily\small} % Spacing and font options for figures
	{\thecontentslabel\hspace*{1em}} % Formatting of numbered sections of this type
	{} % Formatting of numberless sections of this type
	{\ \titlerule*[.5pc]{.}\;\thecontentspage} % Formatting of the filler to the right of the heading and the page number

% Table text styling
\titlecontents{table}
	[1.25cm] % Left indentation
	{\addvspace{1pt}\sffamily\small} % Spacing and font options for tables
	{\thecontentslabel\hspace*{1em}} % Formatting of numbered sections of this type
	{} % Formatting of numberless sections of this type
	{\ \titlerule*[.5pc]{.}\;\thecontentspage} % Formatting of the filler to the right of the heading and the page number

%----------------------------------------------------------------------------------------
%	MINI TABLE OF CONTENTS IN PART HEADS
%----------------------------------------------------------------------------------------

% Chapter text styling
\titlecontents{lchapter}
	[0em] % Left indentation
	{\addvspace{15pt}\large\sffamily\bfseries} % Spacing and font options for chapters
	{\color{uMainColor}\contentslabel[\Large\thecontentslabel]{1.25cm}\color{uMainColor}} % Chapter number
	{}  
	{\color{uMainColor}\normalsize\sffamily\bfseries\;\titlerule*[.5pc]{.}\;\thecontentspage} % Page number

% Section text styling
\titlecontents{lsection}
	[0em] % Left indentation
	{\sffamily\small} % Spacing and font options for sections
	{\contentslabel[\thecontentslabel]{1.25cm}} % Section number
	{}
	{}

% Subsection text styling (note these aren't shown by default, display them by searchings this file for tgcdepth and reading the commented text)
%\titlecontents{lsubsection}
%	[.5em] % Left indentation
%	{\sffamily\footnotesize} % Spacing and font options for subsections
%	{\contentslabel[\thecontentslabel]{1.25cm}}
%	{}
%	{}


\usepackage[]{ccicons}


%----------------------------------------------------------------------------------------
%	THEOREM STYLES
%----------------------------------------------------------------------------------------

\usepackage{amsmath,amsfonts,amssymb,amsthm} % For math equations, theorems, symbols, etc

\newcommand{\intoo}[2]{\mathopen{]}#1\,;#2\mathclose{[}}
\newcommand{\ud}{\mathop{\mathrm{{}d}}\mathopen{}}
\newcommand{\intff}[2]{\mathopen{[}#1\,;#2\mathclose{]}}
\renewcommand{\qedsymbol}{$\blacksquare$}

\ifcsname chapter \endcsname   % verify if command name exists in the CLASS (Book, report, article ...)
	\newtheorem{notation}{Notation}[chapter]
 \else
	\newtheorem{notation}{Notation}[section]
\fi


% Boxed/framed environments
\newtheoremstyle{ocrenumbox}% Theorem style name
{0pt}% Space above
{0pt}% Space below
{\normalfont}% Body font
{}% Indent amount
{\small\bf\sffamily\color{uMainColor}}% Theorem head font
{\;}% Punctuation after theorem head
{0.25em}% Space after theorem head
{\small\sffamily\color{uMainColor}\thmname{#1}\nobreakspace\thmnumber{\@ifnotempty{#1}{}\@upn{#2}}% Theorem text (e.g. Theorem 2.1)
\thmnote{\nobreakspace\the\thm@notefont\sffamily\bfseries\color{black}---\nobreakspace#3.}} % Optional theorem note

\newtheoremstyle{ocrenumbox}% Theorem style name
{0pt}% Space above
{0pt}% Space below
{\normalfont}% Body font
{}% Indent amount
{\small\bf\sffamily\color{uMainColor}}% Theorem head font
{\;}% Punctuation after theorem head
{0.25em}% Space after theorem head
{\small\sffamily\color{uMainColor}\thmname{#1}\nobreakspace\thmnumber{\@ifnotempty{#1}{}\@upn{#2}}% Theorem text (e.g. Theorem 2.1)
\thmnote{\nobreakspace\the\thm@notefont\sffamily\bfseries\color{black}---\nobreakspace#3.}} % Optional theorem note

\newtheoremstyle{blacknumex}% Theorem style name
{5pt}% Space above
{5pt}% Space below
{\normalfont}% Body font
{} % Indent amount
{\small\bf\sffamily}% Theorem head font
{\;}% Punctuation after theorem head
{0.25em}% Space after theorem head
{\small\sffamily{\tiny\ensuremath{\blacksquare}}\nobreakspace\thmname{#1}\nobreakspace\thmnumber{\@ifnotempty{#1}{}\@upn{#2}}% Theorem text (e.g. Theorem 2.1)
\thmnote{\nobreakspace\the\thm@notefont\sffamily\bfseries---\nobreakspace#3.}}% Optional theorem note


\newtheoremstyle{blacknumboxS}% name of the style to be used
{5pt}% measure of space to leave above the theorem. E.g.: 3pt
{5pt}% measure of space to leave below the theorem. E.g.: 3pt
{\normalfont}% name of font to use in the body of the theorem
{}% measure of space to indent
{\small\bf\sffamily}% name of head font
{}% punctuation between head and body
{ }% space after theorem head; " " = normal interword space
{\color{black}\faEye~  \color{uMainColor}\small\sffamily\thmnote{#3 : }}

% Non-boxed/non-framed environments
\newtheoremstyle{ocrenum}% Theorem style name
{5pt}% Space above
{5pt}% Space below
{\normalfont}% Body font
{}% Indent amount
{\small\bf\sffamily\color{uMainColor}}% Theorem head font
{\;}% Punctuation after theorem head
{0.25em}% Space after theorem head
{\small\sffamily\color{uMainColor}\thmname{#1}\nobreakspace\thmnumber{\@ifnotempty{#1}{}\@upn{#2}}% Theorem text (e.g. Theorem 2.1)
\thmnote{\nobreakspace\the\thm@notefont\sffamily\bfseries\color{black}---\nobreakspace#3.}} % Optional theorem note
\makeatother

% Defines the theorem text style for each type of theorem to one of the three styles above



\ifcsname chapter \endcsname

\newcounter{dummy} 

\numberwithin{dummy}{section}
	\theoremstyle{ocrenumbox}
		\newtheorem{theoremeT}[dummy]{Proposition}
		\newtheorem{problem}{Problem}[chapter]
		\newtheorem{exerciseT}{Outillage}[chapter]
	\theoremstyle{blacknumex}
		\newtheorem{exampleT}{Example}[chapter]
	\theoremstyle{blacknumbox}
		\newtheorem{vocabulary}{Vocabulary}[chapter]
		\newtheorem{definitionT}{Definition}[section]
		\newtheorem{corollaryT}[dummy]{Concept}
	\theoremstyle{blacknumboxS}
		\newtheorem{remarqueST}[]{}
		\newtheorem{remarqueT}[dummy]{Remarque}
	\theoremstyle{ocrenum}
	\newtheorem{proposition}[dummy]{Proposition}
  
 \else
 
\newcounter{dummy} 
\numberwithin{dummy}{section}

	\theoremstyle{ocrenumbox}
		\newtheorem{theoremeT}[dummy]{Proposition}
		%\newtheorem{problem}{Problem}[section]
		\newtheorem{exerciseT}{Outillage}[section]
	\theoremstyle{blacknumex}
		\newtheorem{exampleT}{Example}[section]
	\theoremstyle{blacknumbox}
		\newtheorem{vocabulary}{Vocabulary}[section]
		\newtheorem{definitionT}{Definition}[section]
		\newtheorem{corollaryT}[dummy]{concept}
	\theoremstyle{blacknumboxS}
		\newtheorem{remarqueST}[]{}
		\newtheorem{remarqueT}[dummy]{Remarque}
	\theoremstyle{ocrenum}
		\newtheorem{proposition}[dummy]{Proposition}
\fi



%----------------------------------------------------------------------------------------
%	DEFINITION OF COLORED BOXES
%----------------------------------------------------------------------------------------

\RequirePackage[framemethod=default]{mdframed} % Required for creating the theorem, definition, exercise and corollary boxes

% Theorem box
\newmdenv[skipabove=7pt,
skipbelow=7pt,
backgroundcolor=black!5,
linecolor=uMainColor,
innerleftmargin=5pt,
innerrightmargin=5pt,
innertopmargin=5pt,
leftmargin=0cm,
rightmargin=0cm,
innerbottommargin=5pt]{tBox}

% Exercise box	  
\newmdenv[skipabove=7pt,
skipbelow=7pt,
rightline=false,
leftline=true,
topline=false,
bottomline=false,
backgroundcolor=uMainColor!10,
linecolor=uMainColor,
innerleftmargin=5pt,
innerrightmargin=5pt,
innertopmargin=5pt,
innerbottommargin=5pt,
leftmargin=0cm,
rightmargin=0cm,
linewidth=4pt]{eBox}	

% Definition box
\newmdenv[skipabove=7pt,
skipbelow=7pt,
rightline=false,
leftline=true,
topline=false,
bottomline=false,
linecolor=uMainColor,
innerleftmargin=5pt,
innerrightmargin=5pt,
innertopmargin=0pt,
leftmargin=0cm,
rightmargin=0cm,
linewidth=4pt,
innerbottommargin=0pt]{dBox}	

% Corollary box
\newmdenv[skipabove=7pt,
skipbelow=7pt,
rightline=false,
leftline=true,
topline=false,
bottomline=false,
linecolor=gray,
backgroundcolor=black!5,
innerleftmargin=5pt,
innerrightmargin=5pt,
innertopmargin=5pt,
leftmargin=0cm,
rightmargin=0cm,
linewidth=4pt,
innerbottommargin=5pt]{cBox}

% Creates an environment for each type of theorem and assigns it a theorem text style from the "Theorem Styles" section above and a colored box from above
%\newenvironment{theorem}{\begin{tBox}\begin{theoremeT}}{\end{theoremeT}\end{tBox}}
\newenvironment{exercise}{\begin{eBox}\begin{exerciseT}}{\hfill{\color{uMainColor}\tiny\ensuremath{\blacksquare}}\end{exerciseT}\end{eBox}}				  
%\newenvironment{definition}{\begin{dBox}\begin{definitionT}}{\end{definitionT}\end{dBox}}	
%\newenvironment{example}{\begin{exampleT}}{\hfill{\tiny\ensuremath{\blacksquare}}\end{exampleT}}		
%\newenvironment{corollary}{\begin{cBox}\begin{corollaryT}}{\end{corollaryT}\end{cBox}}	
\newenvironment{nota}{\begin{cBox}\begin{remarqueT}}{\end{remarqueT}\end{cBox}}	

%----------------------------------------------------------------------------------------
%	REMARK ENVIRONMENT
%----------------------------------------------------------------------------------------

\newenvironment{remark}{\par\vspace{10pt}\small % Vertical white space above the remark and smaller font size
\begin{list}{}{
\leftmargin=35pt % Indentation on the left
\rightmargin=25pt}\item\ignorespaces % Indentation on the right
\makebox[-2.5pt]{\begin{tikzpicture}[overlay]
\node[draw=uMainColor!60,line width=1pt,circle,fill=uMainColor!25,font=\sffamily\bfseries,inner sep=2pt,outer sep=0pt] at (-15pt,0pt){\textcolor{uMainColor}{i}};\end{tikzpicture}} % Orange R in a circle
\advance\baselineskip -1pt}{\end{list}\vskip5pt} % Tighter line spacing and white space after remark

%----------------------------------------------------------------------------------------
%	SECTION NUMBERING IN THE MARGIN
%----------------------------------------------------------------------------------------

\makeatletter
% \renewcommand{\@seccntformat}[1]{\textcolor{uMainColor}{\csname the#1\endcsname}\quad}   

%\makeatother

%----------------------------------------------------------------------------------------
%	PART HEADINGS
%----------------------------------------------------------------------------------------

% Numbered part in the table of contents
\newcommand{\@mypartnumtocformat}[2]{%
	\setlength\fboxsep{0pt}%
	\noindent\colorbox{uMainColor!20}{\strut\parbox[c][.7cm]{\ecart}{\color{uMainColor!70}\Large\sffamily\bfseries\centering#1}}\hskip\esp\colorbox{uMainColor!40}{\strut\parbox[c][.7cm]{\linewidth-\ecart-\esp}{\Large\sffamily\centering#2}}%
}

% Unnumbered part in the table of contents
\newcommand{\@myparttocformat}[1]{%
	\setlength\fboxsep{0pt}%
	\noindent\colorbox{uMainColor!40}{\strut\parbox[c][.7cm]{\linewidth}{\Large\sffamily\centering#1}}%
}

\newlength\esp
\setlength\esp{4pt}
\newlength\ecart
\setlength\ecart{1.2cm-\esp}
\newcommand{\thepartimage}{}%
\newcommand{\partimage}[1]{\renewcommand{\thepartimage}{#1}}%
\def\@part[#1]#2{%
\ifnum \c@secnumdepth >-2\relax%
\refstepcounter{part}%
\addcontentsline{toc}{part}{\texorpdfstring{\protect\@mypartnumtocformat{\thepart}{#1}}{\partname~\thepart\ ---\ #1}}
\else%
\addcontentsline{toc}{part}{\texorpdfstring{\protect\@myparttocformat{#1}}{#1}}%
\fi%
\startcontents%
\markboth{}{}%
{\thispagestyle{empty}%
\begin{tikzpicture}[remember picture,overlay]%
\node at (current page.north west){\begin{tikzpicture}[remember picture,overlay]%	
\fill[uMainColor!20](0cm,0cm) rectangle (\paperwidth,-\paperheight);
\node[anchor=north] at (4cm,-3.25cm){\color{uMainColor!40}\fontsize{220}{100}\sffamily\bfseries\thepart}; 
\node[anchor=south east] at (\paperwidth-1cm,-\paperheight+1cm){\parbox[t][][t]{8.5cm}{
\printcontents{l}{0}{\setcounter{tocdepth}{1}}% The depth to which the Part mini table of contents displays headings; 0 for chapters only, 1 for chapters and sections and 2 for chapters, sections and subsections
}};
\node[anchor=north east] at (\paperwidth-1.5cm,-3.25cm){\parbox[t][][t]{15cm}{\strut\raggedleft\color{white}\fontsize{30}{30}\sffamily\bfseries#2}};
\end{tikzpicture}};
\end{tikzpicture}}%
\@endpart}
\def\@spart#1{%
\startcontents%
\phantomsection
{\thispagestyle{empty}%
\begin{tikzpicture}[remember picture,overlay]%
\node at (current page.north west){\begin{tikzpicture}[remember picture,overlay]%	
\fill[uMainColor!20](0cm,0cm) rectangle (\paperwidth,-\paperheight);
\node[anchor=north east] at (\paperwidth-1.5cm,-3.25cm){\parbox[t][][t]{15cm}{\strut\raggedleft\color{white}\fontsize{30}{30}\sffamily\bfseries#1}};
\end{tikzpicture}};
\end{tikzpicture}}
\addcontentsline{toc}{part}{\texorpdfstring{%
\setlength\fboxsep{0pt}%
\noindent\protect\colorbox{uMainColor!40}{\strut\protect\parbox[c][.7cm]{\linewidth}{\Large\sffamily\protect\centering #1\quad\mbox{}}}}{#1}}%
\@endpart}
\def\@endpart{\vfil\newpage
\if@twoside
\if@openright
\null
\thispagestyle{empty}%
\newpage
\fi
\fi
\if@tempswa
\twocolumn
\fi}

%----------------------------------------------------------------------------------------
%	CHAPTER HEADINGS
%-------------------------------------------------%--------------------------------------

%
%
%% A switch to conditionally include a picture, implemented by Christian Hupfer
\newif\ifusechapterimage
\usechapterimagetrue
\newcommand{\thechapterimage}{}%
\newcommand{\chapterimage}[1]{\ifusechapterimage\renewcommand{\thechapterimage}{#1}\fi}%
\newcommand{\autodot}{.}

\def\@makechapterhead#1{%
{\parindent \z@ \raggedright \normalfont
\ifnum \c@secnumdepth >\m@ne
\if@mainmatter
\begin{tikzpicture}[remember picture,overlay]
	
\node at (current page.north west)
			{\begin{tikzpicture}[remember picture,overlay]
			\node[anchor=north west,inner sep=0pt] at (0,0) {\ifusechapterimage\includegraphics[width=\paperwidth]{\thechapterimage}\fi};
			\draw[anchor=west] (\Gm@lmargin,-9cm) node [line width=2pt,rounded corners=15pt,draw=uMainColor,fill=white,fill opacity=1,inner sep=15pt]{\strut\makebox[22cm]{}};
			
			\draw[anchor=west] (\Gm@lmargin+.3cm,-9cm) node {\fontsize{20}{30}\selectfont\sffamily\bfseries\color{uMainColor}~#1\strut};
			%{\Huge\sffamily\bfseries\color{uMainColor}\thechapter\autodot~#1\strut};
			\node at (current page.north east)  [xshift=-3.95cm, yshift=-4.05cm, text opacity=1]  {{\color{white}\centering\fontsize{200}{30}\selectfont \bfseries\sffamily\thechapter\strut} };
			\node at (current page.north east)  [xshift=-4cm, yshift=-4cm, text opacity=1]  {{\color{uMainColor}\centering\fontsize{200}{30}\selectfont \bfseries\sffamily\thechapter\strut} };
			
			\end{tikzpicture}};
	\end{tikzpicture}

\else

\begin{tikzpicture}[remember picture,overlay]
\node at (current page.north west)
{\begin{tikzpicture}[remember picture,overlay]
\node[anchor=north west,inner sep=0pt] at (0,0) {\ifusechapterimage\includegraphics[width=\paperwidth]{\thechapterimage}\fi};
\draw[anchor=west] (\Gm@lmargin,-9cm) node [line width=2pt,rounded corners=15pt,draw=uMainColor,fill=white,fill opacity=1,inner sep=15pt]{\strut\makebox[22cm]{}};

\draw[anchor=east] (\Gm@lmargin,-9cm) node [line width=2pt,rounded corners=15pt,draw=uMainColor,fill=white,fill opacity=1,inner sep=15pt]{\strut\makebox[22cm]{}};
\draw[anchor=west] (\Gm@lmargin+.3cm,-9cm) node {\huge\sffamily\bfseries\color{uMainColor}#1\strut};

\end{tikzpicture}};
 
\end{tikzpicture}
\fi\fi\par\vspace*{270\p@}}}

\makeatother
%-------------------------------------------

%\def\@makeschapterhead#1{%
%\begin{tikzpicture}[remember picture,overlay]
%\node at (current page.north west)
%{\begin{tikzpicture}[remember picture,overlay]
%\node[anchor=north west,inner sep=0pt] at (0,0) {\ifusechapterimage\includegraphics[width=\paperwidth]{\thechapterimage}\fi};
%\draw[anchor=west] (\Gm@lmargin,-9cm) node [line width=2pt,rounded corners=15pt,draw=uMainColor,fill=white, fill opacity=0.8,inner sep=15pt]{\strut\makebox[22cm]{}};
%\draw[anchor=west] (\Gm@lmargin+.3cm,-9cm) node {\huge\sffamily\bfseries\color{uMainColor}#1\strut};
%\end{tikzpicture}};
%\end{tikzpicture}
%\par\vspace*{270\p@}}




%----------------------------------------------------------------------------------------
%	LINKS
%----------------------------------------------------------------------------------------


\usepackage{bookmark}
\bookmarksetup{
open,
numbered,
addtohook={%
\ifnum\bookmarkget{level}=0 % chapter
\bookmarksetup{bold}%
\fi
\ifnum\bookmarkget{level}=-1 % part
\bookmarksetup{color=uMainColor,bold}%
\fi
}
}

%----------------------------------------------------------------------------------------
%	FANCY CHAPTER Header
%----------------------------------------------------------------------------------------





%\definecolor{uMainColor}{RGB}{255,128,0}   

%-------------------------------------------------
%	BIBLIOGRAPHY
%-------------------------------------------------

\addbibresource{Latex/Sources/commons/bibliography.bib} % BibTeX bibliography file
\defbibheading{bibempty}{}

%-------------------------------------------------
%	COLORS & BORDERS
%-------------------------------------------------
\setlength{\fboxrule}{0.75pt} % Width of the border around the abstract
%\definecolor{uMainColor}{RGB}{255,128,0} % Color of the article title and section
\definecolor{color2}{RGB}{220,220,220} % Color of the boxes behind the abstract and headings

%-------------------------------------------------
% ITEMS DEFINITION
%-------------------------------------------------

\setlist [itemize,1]{label=\color{uMainColor}\faCaretRight }

\begin{document}


\addtocontents{toc}{\protect\hypertarget{toc}{}}

%-------------------------------------------------
%	ARTICLE INFORMATION
%-------------------------------------------------

\Abstract {\uabstract%===========================
% COURS "INTRO CYBERDEF"
% Abstract général des articles
%===========================


Il fait partie du cours introductif aux fondamentaux de la sécurité des systèmes d'information vue sous deux prismes quelques fois opposés dans la littérature : la gouvernance et la gestion opérationnelle de la sécurité.
Le cours est constitué d'un ensemble de notes de synthèse indépendantes compilées en un document unique, mais édité par chapitre dans le cadre de ce cours.\\
C\edoc ne constitue pas à lui seul le référentiel du cours CYBERDEF101 (SEC101 du Cnam). Il compile des notes de cours mises à disposition de l'auditeur comme support pédagogique partiel à ce cours introductif à la cyberdéfense d'entreprise. }
\JournalInfo{\uJournalInfo} % Journal information
\Archive{Notes de cours éditées le  \DTMnow} % Additional notes (e.g. copyright, DOI, review/research article)
\PaperTitle{\utitle} % Article title
\Authors{\uauthor\textsuperscript{1,}\textsuperscript{2}*} % Authors
\affiliation{\textsuperscript{1}\textit{\uproa}} % Author affiliation
\affiliation{\textsuperscript{2}\textit{\uprob}} % Author affiliation
\affiliation{*\textbf{email}: \umaila\ -- \umailb} % Corresponding author
\Keywords{\ukeywords} % Keywords 
\newcommand{\keywordname}{Mots clefs} % Defines the keywords heading name

%-------------------------------------------------
%	MAKE TITLE
%-------------------------------------------------

%\maketitle

%\begin{tikzpicture}[remember picture,overlay]\node at (current page.north west)[xshift=4cm, yshift=-4cm, text opacity=1.0]{\includegraphics[width=0.1\paperwidth]{\utemplatepath/commons.inc/CommonsPictures/\utemplatetheme.theme/shield-20.pdf}}; 
%\end{tikzpicture} 
%
%
%\AddToShipoutPicture{%
%\AtPageCenter{%
%\begin{tikzpicture}[remember picture,overlay]
%\node [rotate=60,scale=5,text opacity=0.1] at (current page.center) { \textit{\color{uAlertColor}\INFODistrib} };
%\end{tikzpicture} 
%%\begin{tikzpicture}[remember picture,overlay]
%%\node [rotate=60,scale=10,text opacity=0.3]
%%at (current page.center) {moi};
%%\end{tikzpicture}
%}%
%}


%-------------------------------------------------	
%\begin{center}
%	\textcolor{uMainColor}{\textbf{\INFODistrib}} % Distribution information
%\end{center}
%
%\INFOGitHub % Link to download file and copyright
%%-------------------------------------------------

%%-------------------------------------------------------------
%               FR CYBERDEF SECOPS COURSE
%
%                 DEFINITION DES Sources/EDU/acronymes

%                      $File :  Sources/EDU/acronymes.tex
%                             2020 eduf@ction
%-------------------------------------------------------------
% Require : glossaries package
%-------------------------------------------------------------
%\newacronym{XXX}{XXX}{XXX}
%-----------------------------------------------------------

%  Organismes

\newacronym{aCNIL}{CNIL}{Commission National Informatique et Liberté}
\newacronym{aOTAN}{OTAN}{Organisation du Traité de l'Atlantique Nord}
\newacronym{aANSSI}{ANSSI}{Agence Nationale de la Sécurité des systèmes d'information}
\newacronym{aCERT}{CERT}{Computer Emergency Response Team}
\newacronym{aAFNOR}{AFNOR}{Association Française de Normalisation}
\newacronym{aSIEM}{SIEM}{System Incident and Event Management}
\newacronym{aSIM}{SIM}{System Incident Management}
\newacronym{aSEM}{SEM}{System Event Management}
\newacronym{aIRP}{IRP}{Incident Response Plan, Plan de Réponse à Incident}
\newacronym{aISO}{ISO}{International Organization for Standardization}
\newacronym{aDICT}{DICT}{Disponibilité, Intégrité, Confidentialité, Traçabilité}
\newacronym{aPDCA}{PDCA}{Plan, Do, Check, Act (Roue de Deming)}
\newacronym{aGDPR}{GDPR}{General Data Protection Regulation ou RGPD Règlement général sur la protection des données}

%-----------------------------------------------------------

% ROLES ET METIERS

\newacronym{aOS}{OS}{Officier de sécurité}
\newacronym{aRSSI}{RSSI}{Responsable de la sécurité des systèmes d'information}
\newacronym{aDPO}{DPO}{Data Protection Officier}
\newacronym{aCIL}{CIL}{ Correspondant Informatique et Liberté}
\newacronym{aSSI}{SSI}{Sécurité des Systèmes d’information}
\newacronym{aGRC}{GRC}{Gouvernance Risques et Conformité}
\newacronym{aDTF}{DTF}{Defend The Flag}
\newacronym{aCTF}{CTF}{Capture The Flag}
\newacronym{aOIV}{OIV}{Opérateur d'Infrastructures Vitales}
\newacronym{aOSE}{OSE}{Opérateur de Services Essentiels}
\newacronym{aLPM}{LPM}{Loi de Programmation Militaire}
\newacronym{aEBIOS}{EBIOS}{Expression des Besoins et Identification des Objectifs de Sécurité}
\newacronym{aDLL}{DLL}{Dynamic Loads Library}
\newacronym{aACID}{ACID}{Authentification, Confidentialité, Intégrité, Disponibilité}

%-------------------------------------------------------
% Sources/EDU/acronymes DATABASE XLS Generation (cf. RESSOURCES DIR)
%-------------------------------------------------------
% DO NOT MODIFY DIRECTLY, USE XLS File (BEGIN)

\newacronym{a2FA}{2FA}{\textit{Two-Factor Authentication} : Authentification à deux facteurs}
\newacronym{aAAA}{AAA}{\textit{Authentication, Authorization, Accounting} : Authentification, autorisation, comptabilisation}
\newacronym{aAC}{AC}{\textit{Access Control} : Accès conditionnel}
\newacronym{aAES}{AES}{\textit{Advanced Encryption Standard} : Algorithme de chiffrement avancé}
\newacronym{aAPT}{APT}{\textit{Advanced Persistent Threat Group} : Groupe de cyber-espionnage avancé}
\newacronym{aASLR}{ASLR}{\textit{Address Space Layout Randomization} : Randomisation de l'espace d'adressage}
\newacronym{aAV}{AV}{\textit{Antivirus Software} : Logiciel antivirus}
\newacronym{aBOT}{BOT}{\textit{Robot (Malicious)} : Robot (malveillant)}
\newacronym{aCC}{C\&C}{\textit{Command And Control} : Commande et contrôle}
\newacronym{aCSRF}{CSRF}{\textit{Cross-Site Request Forgery Attack} : Attaque de confiance de site à site}
\newacronym{aDDOS}{DDOS}{\textit{Distributed Denial Of Service Attack} : Attaque de déni de service distribué}
\newacronym{aDES}{DES}{\textit{Data Encryption Standard} : Algorithme de chiffrement à clé de 64 bits}
\newacronym{aDLP}{DLP}{\textit{Data Loss Prevention} : Gestion de la protection des données sensibles}
\newacronym{aEDR}{EDR}{\textit{Endpoint Detection And Response} : Détection et réponse aux comportements anormaux d'un terminal}
\newacronym{aEK}{EK}{\textit{Exploit Kit} : Kit d'exploit}
\newacronym{aFTP}{FTP}{\textit{File Transfer Protocol} : Protocole de transfert de fichier}
\newacronym{aHTTPS}{HTTPS}{\textit{Protocol For Securing Communications On The Internet} : Protocole de sécurisation HTTP}
\newacronym{aIAM}{IAM}{\textit{Identity And Access Management} : Identité et Gestion des Accés}
\newacronym{aIDS}{IDS}{\textit{Intrusion Detection System} : Système de détection d'intrusion}
\newacronym{aIIS}{IIS}{\textit{Internet Information Services Server}}
\newacronym{aIPS}{IPS}{\textit{Intrusion Prevention System} : Système de prévention d'incident}
\newacronym{aISA}{ISA}{\textit{Internet Security And Acceleration Server} : Serveur de sécurité et d'accélération Internet (Microsoft)}
\newacronym{aLDAP}{LDAP}{\textit{Lightweight Directory Access Protocol} : Protocole de gestion d'annuaire léger}
\newacronym{aMD5}{MD5}{\textit{Message-Digest Algorithm 5}}
\newacronym{aMFA}{MFA}{\textit{Multi-Factor Authentication} : Authentification multi-facteurs}
\newacronym{aNAC}{NAC}{\textit{Network Access Control}}
\newacronym{aOAuth}{OAuth}{\textit{Open Authorization Protocol}}
\newacronym{aPAM}{PAM}{\textit{Privilege Access Management}}
\newacronym{aPKI}{PKI}{\textit{Public Key Infrastructure} : Infrastructure de clefs publiques}
\newacronym{aRSA}{RSA}{\textit{Rivest-Shamir-Adleman Encryption Algorithm}}
\newacronym{aSAML}{SAML}{\textit{Security Assertion Markup Language}}
\newacronym{aSMS}{SMS}{\textit{Short Message Service}}
\newacronym{aSQLI}{SQLI}{\textit{Sql Injection} : Injection de script de commande SQL}
\newacronym{aSSH}{SSH}{\textit{Secure Shell Protocol}}
\newacronym{aSSL}{SSL}{\textit{Protocol For Securing Communications On The Internet} : Protocole de sécurisation des communications sur Internet}
\newacronym{aTKIP}{TKIP}{\textit{Temporal Key Integrity Protocol}}
\newacronym{aTLS}{TLS}{\textit{Transport Layer Security}}
\newacronym{aURL}{URL}{\textit{Uniform Resource Locator}}
\newacronym{aVPN}{VPN}{\textit{Virtual Private Network} : Réseau privé virtuel}
\newacronym{aWAF}{WAF}{\textit{Web Application Firewall}}
\newacronym{aWPA}{WPA}{\textit{Wifi Protected Access}}
\newacronym{aXSS}{XSS}{\textit{Cross-Site Scripting Attack}}

% DO NOT MODIFY (END)
%-------------------------------------------------------





%-------------------------------------------------
\newpage
\ubody % MAIN BODY defined external
%-------------------------------------------------

	\printbibliography
 \printglossary[type=\acronymtype]
%	\printglossary
%-------------------------------------------------	
\newpage
\Ucontribute % Contribution TEXT
%-------------------------------------------------


%-------------------------------------------------
	\newpage
	\tableofcontents  % TOC
	\iftotalfigures
  			\listoffigures
	\fi
%-------------------------------------------------

\end{document}
}
	\else
	\providecommand{\umainload}{%-------------------------------------------------------------
%               FR CYBERDEF SECOPS COURSE
%                          ARTICLE MAIN FILE
%                             2020 eduf@ction
%-------------------------------------------------------------

\documentclass[10pt,fleqn,twoside]{\utemplatepath/LessonModel/Article/edxstyle} %edxstyle

\usepackage{beamerarticle}

\usepackage{ae}
\usepackage[french]{babel}
\usepackage[utf8]{inputenc}
\usepackage[T1]{fontenc}
\usepackage{lmodern}




\usepackage{titletoc} % Required for manipulating the table of contents

\input{\utemplatepath/commons.inc/CommonsPictures/\utemplatetheme.theme/colors.inc}

%-------------------------------------------------------------
%               FR CYBERDEF SECOPS COURSE
%                                  EDX PACKAGES
%.    Commons packages for BOOK, NOTES (.doc) 
%                            and BEAMER (.prz)
%                              2020 eduf@ction
%-------------------------------------------------------------

%\usepackage[hyphens,spaces,obeyspaces]{url}
%\usepackage[hyphens,spaces,obeyspaces]{url}



\PassOptionsToPackage{hyphens,spaces,obeyspaces}{url}

\usepackage[utf8]{inputenc}
\usepackage[french]{babel}
\usepackage[T1]{fontenc}
%\usepackage{markdown}



\usepackage{ccicons}

\usepackage{tikz} % Required for drawing custom shapes
\usetikzlibrary{tikzmark} 
\usetikzlibrary{mindmap,trees, backgrounds}


\usepackage{verbatim}
\usepackage{wrapfig}
\usepackage{xpatch}

\usepackage{hologo}

\usepackage[figure]{totalcount}

\usepackage{fontawesome}

\usepackage[skins]{tcolorbox} % don't forget SKINS option

\usepackage{listings}
\usepackage{upquote}

\lstset{
upquote=true,columns=flexible, basicstyle=\ttfamily,
language=HTML, 
frameround=tttt,
commentstyle=\color{gray},
identifierstyle=\color{blue},
keywordstyle=\color{uMainColor}\bfseries,
xleftmargin=2em,
xrightmargin=2em,
aboveskip=\topsep,
belowskip=\topsep, 
frame=single,
rulecolor=\color{uMainColor},
backgroundcolor=\color{uMainColor!5},
breaklines,
breakindent=1.5em,
showspaces=false,
showstringspaces=false,
showtabs=false,
xleftmargin=2em,
xrightmargin=2em,
aboveskip=1em,
belowskip=1em,
}



\usepackage{graphicx} % Required for including pictures
%\graphicspath{{Pictures/}} % Specifies the directory where pictures are stored

\usepackage{booktabs} % Required for nicer horizontal rules in tables

\usepackage{setspace}
\setstretch{1,2}

\usepackage{csquotes}

\usepackage{fancyhdr}

\usepackage{makeidx} 

\usepackage{xspace} 

\usepackage{floatflt} 

\usepackage{hyperref}
\usepackage{cleveref}

%pdfdisplaydoctitle={\utitle}, Attention le titre ne doit pas contenir de saut de ligne 

\hypersetup{
pdftitle={\uCoursetittle},
pdfauthor={\uauthor},
hidelinks,
backref=true,
pagebackref=true,
hyperindex=true,
pdfcenterwindow = true,
colorlinks=true,
breaklinks=true,
bookmarks=true,
allcolors=uMainColor,
pdfsubject = {\udescription},
pdfcreator = {MACTEX on TEXPAD},
bookmarksopen=true}


% urlcolor=cnam,
% runcolor=uMainColor,
%\pdfstringdefDisableCommands{%
%  \def\\{}%
%  \def\texttt#1{<#1>}%
%}

\usepackage[style=numeric,citestyle=numeric,sorting=nyt,sortcites=true,autopunct=true,hyperref=true,abbreviate=false,backref=true,backend=biber]{biblatex}

%\usepackage{cleveref} old position

\usepackage{multicol}

%--------------------------------------------------------
% pour la génération du QRCODE en page de garde d'article
\usepackage{qrcode} 

%--------------------------------------------------------
% POur la date d'impression en premiere page d'article
\usepackage[useregional]{datetime2}

%--------------------------------------------------------
% Gestion des glossaires et Sources/EDU/acronymes
% Un fichier acroymes.tex centralisé
\usepackage[acronym]{glossaries}

%\makeglossaries
\makenoidxglossaries 

\loadglsentries{Latex/Sources/commons/acronymes.tex}
\loadglsentries{Latex/Sources/commons/glossaire.tex}


%\renewcommand{\bibname}{R\'ef\'erences}

%--------------------------------------------------------
% Pour exclure des "environnement", ici techworkbox par exemple, travaux pour étudiants, non affiché dans le BOOK.
\usepackage{comment}


%--------------------------------------------------------
% pour filigramme sur les pages
\usepackage{eso-pic}

%--------------------------------------------------------
% for \gappto{\UrlBreaks}{\UrlOrds}
% rather than 
%\makeatletter
% \g@addto@macro{\UrlBreaks}{\UrlOrds}
% \makeatother

\usepackage{etoolbox}

\usepackage{breakurl}


\newcommand{\ucite}[1]{\cite{#1}}
%\newcommand{\ucite}[1]{REF BIB #1}
\newcommand{\ugls}[1]{\gls{#1}}
%\newcommand{\ugls}[1]{GLOSSAIRE (#1)}
%\usepackage[Rejne]{fncychap}

%\usepackage{pythontex} % need shell escape

%
%\makeatletter
%\DeclareRobustCommand{\afilename}[1]{%
% \begingroup
%  % \lstname seems to change hyphens into \textendash
%  \def\textendash{-}%
%  \afilename@parse{#1}%
%  \edef\afilename@base{\detokenize\expandafter{\afilename@base}}%
%  \texttt{\afilename@base.\afilename@ext}%
% \endgroup
%}
%\makeatother

%-------------------------------------------------------------
%               FR CYBERDEF SECOPS COURSE
%                                  EDX PACKAGES
%.    Commons COMMANDS for BOOK, NOTES (.doc) 
%                            and BEAMER (.prz)
%                              2020 eduf@ction
%-------------------------------------------------------------

\newcommand{\ucite}[1]{\cite{#1}}
%\newcommand{\ucite}[1]{REF BIB #1}
\newcommand{\ugls}[1]{\gls{#1}}
%\newcommand{\ugls}[1]{GLOSSAIRE (#1)}
\newcommand{\uac}[1]{\acrshort{#1}}



% \newcommand{\ExtractBeforeLastSlash}[1]{\directlua{extract_before_last_slash("#1")}}
% \newcommand{\ExtractAfterLastSlash}[1]{\directlua{extract_after_last_slash("#1")}}
% \newcommand{\RemoveExtension}[1]{\directlua{remove_extension("#1")}}
% \newcommand{\ExtractExtension}[1]{\directlua{extract_extension("#1")}}


\newcommand{\includepicturesubpath}[1]
{#1}

%{\ExtractBeforeLastSlash{#1}/\upicturethemepath\RemoveExtension{\ExtractAfterLastSlash{#1}}.\ExtractExtension{#1}}
    

%---------------------------------------------------
%										 	AUDIO NOTE 
%---------------------------------------------------

\newcommand\audionote[2]{%
  \note{#2}
  \immediate\write\audioexport{==AudioSlide=============}
  \immediate\write\audioexport{#1}
  \immediate\write\audioexport{\insertframenumber}
  \immediate\write\audioexport{---------------------------------}
  \immediate\write\audioexport{#2}
  \immediate\write\audioexport{---------------------------------}
}%

%---------------------------------------------------
%										 	TEXT STYLE SHORTCUT 
%---------------------------------------------------

\newcommand{\tcb}[1]{\textbf{\textcolor{uMainColor}{#1}}}

\newcommand{\voir}[1]{(Schéma~\ref{#1},~page~\pageref{#1}}

\newcommand{\graphicxextension}{pdf}

\newcommand\myemptypage{
   
    \thispagestyle{empty}
    \addtocounter{page}{-1}
    \newpage
  \null
    %CECI EST UNE PAGE VIDE
       \vfill
    }


\newcommand{\uNA}[2]
{
\newacronym{a#1}{#1}{#2}
}

\newcommand{\coauthor}[4]
{
\noindent (#2) \textbf{#1} - #4 : \textit{#3}
}


\newcommand\UKword[1]{\emph{#1}\xspace}
\newcommand\textUK[1]{\textit{\textbf{#1}}\xspace}
\newcommand\Ucenter[1]{\begin{center}{#1}\end{center}}

\newcommand\head[1]{\textbf{#1}}
\newcommand\tb[1]{\textbf{#1}}
\newcommand{\g}[1]{\og #1 \fg{}}
\newcommand{\uload}[1] {\input{../Latex/#1}}
\newcommand{\ucurcolor}{black}

\newcommand{\mak}[1]{\faicon{\aArrowCircleORight} #1}


%-----------------------------------------------------
%										 	\INFOGitHub 
%---------------------------------------------------


\newcommand{\GITfilename}{https://github.com/edufaction/CYBERDEF101/raw/main/Release/Builder/\jobname.pdf}

%\newcommand{\safeqrcode}[2][]{%
%\qrcode[#1]{\detokenize{#2}} }

\newcommand{\safeqrcode}[2][]{\qrcode[#1]{#2}} 

\newcommand{\INFOGitHub}
{
\begin{center} 
\setstretch{1.5}
 {Vérifiez la disponibilité d'une version plus récente de} \\
 {\ulink{\GITfilename}{\textbf{\jobname.pdf }sur GITHUB CYBERDEF} }   \\ 
{{\huge\ccbyncndeu}}  \\  
{Publication en \href{https://creativecommons.org/licenses/by-nc-nd/2.0/fr/}{Creative Common BY-NC-ND} by eduf@ction}    \\  % Copyright notice
{\safeqrcode[height=4 cm,padding]{\GITfilename}}  
\end{center} 
}

\newcommand{\INFOGitHubS}
{
\noindent{\safeqrcode[height=3 cm]{\GITfilename}} \\
\noindent {\tiny Téléchargez une version à jour} 
}



%-----------------------------------------------------
%										 	\upspicture 
%---------------------------------------------------
\newcommand{\upspicture}[3]
{
\renewcommand{\ucurcolor}{#3}
\resizebox{#2\textwidth}{!}{\input{../#1}}
}


%-----------------------------------------------------
%										 	\uindex 
%---------------------------------------------------


\newcommand{\uindex}[1]
{
#1\index{#1}\xspace
}

\newcommand{\ulindex}[1]{\index{#1}}

%-----------------------------------------------------
%										 	\includer 
%---------------------------------------------------

\newcommand{\includer}[1]
{
\input{../Latex/Sources/EDU/SRC1/Chapters/#1}
}

%-----------------------------------------------------
%										 	\edxdico 
%---------------------------------------------------

\newcommand{\edxdico}[2] {\textbf{#1}:\textit{#2} \index{#1}}

%-----------------------------------------------------
%  command DEF : InTextx
%-----------------------------------------------------
\newcommand\InTexta{no text}
\newcommand\InTextb{no text}
\newcommand\InTextc{no text}
\newcommand\rnc{\renewcommand}


%-----------------------------------------------------
%										 	ENV warningbox
%---------------------------------------------------

\newtcolorbox{warningbox}[2][]
{
  colframe = uMainColor!25,
  colback  = black!05,
  coltitle = uMainColor!20!black,
  fonttitle = \bfseries, 
  title    = #2,
  #1,
}


\newtcolorbox{techworkbox}[2][]{sharp corners,colframe = uMainColor!25,colback  = black!05,coltitle = uMainColor!20!black,fonttitle = \bfseries, title={\textcolor{uMainColor}{\faFile~\small {Sujet de mémoire SECOPS}}} : #2,#1,}


\newtcolorbox{notebox2}[2][]
{
sharp corners, 
colback = uMainColor!5!white, 
colframe = uMainColor!75!black,
fonttitle = \bfseries, 
colbacktitle= uMainColor!85!black,
title=#2,#1
}

\newtcolorbox{notebox}[2][]{colback=uMainColor!5!white,
colframe=uMainColor!75!black,fonttitle=\bfseries,
colbacktitle=uMainColor!85!black,enhanced,
attach boxed title to top right={yshift=+1mm},
title=#2,#1}


\newtcolorbox{toolsboxin}[2][]{sharp corners, colback=uMainColor!0!white,
colframe= black!25,fonttitle=\bfseries,
colbacktitle=black!60,enhanced,
attach boxed title to top left={yshift=+1mm},
title=#2,#1}


\newtcolorbox{techworkin}[2][]{sharp corners, colback=uMainColor!5!white,
colframe= uMainColor!100, fonttitle=\bfseries,
colbacktitle=black!5, enhanced,
 coltitle=black!50,
attach boxed title to top right ={yshift=+1mm},
title=#2,#1}



%-----------------------------------------------------
%										 	ENV ToolsBox
%---------------------------------------------------


\newcommand\tooleditor{No EDITOR in Datas.tex}
\newcommand\toolurl{No URL INSERTED in Datas.tex}
\newcommand\toolclass{No EDITOR in Datas.tex}
\newcommand\toolname{No NAME in Datas.tex}
\newcommand\toolanalyst{No ANALYST in Datas.tex}

\newcommand{\toolsbox}[2]
{
\input{../Latex/Sources/EDU/SRC1/SecComponents/#2.#1/datas.tex} 
\begin{toolsboxin}[sidebyside,righthand width=.7\textwidth]{#1}
\includegraphics[width=0.9\textwidth]{../Latex/Sources/EDU/SRC1/SecComponents/#2.#1/icon.png} \\
\tcblower
\input{../Latex/Sources/EDU/SRC1/SecComponents/#2.#1/description.tex}
\end{toolsboxin}
\noindent \small {\faCogs~Classe : \uppercase{\textbf{#2}}, Site de référence : \textbf{\ulink{\toolurl}{#1}}}
\newline
\small {\faGlobe~Editeur : \textbf{\tooleditor}}
\small {\faEye~Analyste : \textbf{\toolanalyst}}
}

%\renewcommand\toolsclass{SIEM}
%\renewcommand\toolsname{RSA Netwitness}
%\renewcommand\toolseditor{RSA}
%\renewcommand\toolsurl{https://www.rsa.com/fr-fr/products/threat-detection-response}

%-----------------------------------------------------
%										 	ENV  TechWork
%---------------------------------------------------


\newcommand{\techwork}[2]
{
\begin{techworkin} [righthand width=.8\textwidth]{\faCogs~\small {Fiche TECHNO} : #1}#2
\end{techworkin}
}

%-----------------------------------------------------
%										 	\uchap
%---------------------------------------------------
\ifdefined\INTERTITLE
\providecommand{\uchap}[1]
	{
	\begin{LARGE}
	\textbf{\Ucenter{#1 \\ -oOo-}}	
	\end{LARGE}
	}
\else
	\providecommand{\uchap}[1]{}
\fi

%-----------------------------------------------------
%											\uexpand
%---------------------------------------------------
\newcommand {\uexpand}[1]{
\mode<all>\input{../Latex/Sources/EDU/SRC1/Chapters/#1}
}


%-----------------------------------------------------
%											\spicture 
% direct include graphic with Scrivener
%-----------------------------------------------------

\newcommand{\spicture}[3]{
\begin{figure}[!hbtp] %h !hbtp
  \begin{center}
	 \includegraphics[width=0.8\textwidth]{#1}
  \end{center}
\caption{#3 - (slbl-#2)\label{slbl-#2}}
\end{figure}
}

%\StrBefore{#2}{.pdf}

%-----------------------------------------------------
%											\upicture 
%---------------------------------------------------

\ifdefined\PRZMODE

\providecommand{\upicture}[4]{
\framesubtitle{#2}
%\begin{figure}[!h] %h
  \begin{center}
	 \includegraphics[height=#3\textheight, width=#3\textwidth, keepaspectratio ]{\includepicturesubpath{#1.pdf}}
  \end{center}
%  \end{figure}
}

\else

\providecommand{\upicture}[4]{
\begin{figure}[!hbtp] %h !hbtp
  \begin{center}
	 \includegraphics[width=#3\textwidth]{\includepicturesubpath{#1.pdf}}
  \end{center}
\caption{#2 \label{#4}}
\end{figure}
}
	
\fi


\ifdefined\PRZMODE

\providecommand{\uwpicture}[4]{
\framesubtitle{#2}
%\begin{figure}[!h] %h
  \begin{center}
	 \includegraphics[height=0.8\textheight, width=0.8\textwidth, keepaspectratio ]{\includepicturesubpath{#1.pdf}}
  \end{center}
%  \end{figure}
}
\else
\providecommand{\uwpicture}[4]
{
\begin{wrapfigure}{O}{#3\textwidth} % inside / Outside
  \begin{center}
	 \includegraphics[width=#3\textwidth]{\includepicturesubpath{#1.pdf}}
\end{center}
\caption{#2 \label{#4}}
\end{wrapfigure}
}
\fi

%-----------------------------------------------------
%	\picframe for BEAMER (Image included in Beamer Slide)
%---------------------------------------------------

\ifdefined\PRZMODE

\providecommand{\picframe}[4]
{
\begin{frame}
\frametitle<presentation>{#2}
  \begin{center}

   \begin{tikzpicture}[remember picture, overlay]
    \node at (current page.center) {\includegraphics[width=#3\paperwidth, height=#3\paperheight,keepaspectratio]{\includepicturesubpath{#1.pdf}}};
    \end{tikzpicture}

  \end{center}
\end{frame}

}
\else
\providecommand{\picframe}[4]
{
\begin{figure}[!hbtp] % !hbtp ou !h 
\centering 
%{\includegraphics[width=#3\paperwidth, height=#3\paperheight,keepaspectratio]{\includepicturesubpath{#1.pdf}}}
\includegraphics[width=#3\textwidth]{\includepicturesubpath{#1.pdf}}
\caption{#2 \label{#4}}
\end{figure}
}
\fi

%-----------------------------------------------------
%.  for BEAMER (Image included in Beamer Slide)
%	\picframeS Simplified with Name and Label from One Param
%---------------------------------------------------


\ifdefined\PRZMODE

\providecommand{\picframeS}[2]
{
\begin{frame}
\frametitle<presentation>{#2}
  \begin{center}
	 \includegraphics[height=0.95\textheight, width=0.95\textwidth, keepaspectratio ]{\picpath/img-#1.\graphicxextension}
  \end{center}
\end{frame}

}
\else
\providecommand{\picframeS}[2]
{
\begin{figure}[!hbtp] % !hbtp ou !h 
\centering \includegraphics[width=0.9\textwidth]{\picpath/img-#1.\graphicxextension}
\caption{(s) #2 \label{lbl:#1}}
\end{figure}
}

\fi


%{figure}[!h]
%{wrapfigure}{O}{#3\textwidth}
%-----------------------------------------------------
%	FRAME			\wpicframe 
%---------------------------------------------------

\ifdefined\PRZMODE

\providecommand{\wpicframe}[4]
{
\begin{frame}
\frametitle<presentation>{#2}
  \begin{center}
	 \includegraphics[height=0.95\textheight, width=0.95\textwidth, keepaspectratio ]{#1.\graphicxextension}
  \end{center}
\end{frame}

}
\else
\providecommand{\wpicframe}[4]
{
\begin{wrapfigure}{O}{#3\textwidth}
\centering \includegraphics[width=#3\textwidth]{#1.\graphicxextension}
\caption{#2 \label{#4}}
\end{wrapfigure}
}
\fi

%-----------------------------------------------------
%	FRAME										\texframe 
%---------------------------------------------------
\ifdefined\PRZMODE

\providecommand{\texframe}[3]
{
\begin{frame}
\frametitle<presentation>{#1}
\framesubtitle<presentation>{#2}
#3
\end{frame}
}
\else
\providecommand{\texframe}[3]
{
#3
}

\fi

%-----------------------------------------------------
%											\uref 
%---------------------------------------------------
%\newcommand{\uref}[2]{(Voir~#1~\cref{#2} page~\pageref{#2})}

\newcommand{\uref}[2]{(Voir~#1~\ref{#2} page~\pageref{#2})}
%-----------------------------------------------------
%											\ubg
%---------------------------------------------------
\newcommand{\ubg}[1]{
\textbf{\g{\uppercase{#1}}}
}
%-----------------------------------------------------
%											\updfimage 
%---------------------------------------------------

\newcommand{\updfimage}[3]
{
\begin{wrapfigure}{R}{#3\textwidth}
  \begin{center}
	 \resizebox{#3\textwidth}{!}{\includegraphics{\includepicturesubpath{#1.pdf}}}
  \end{center}
\caption{#2}
\end{wrapfigure}
}


%************************PACKAGE***

%\usepackage{tikz}% Required for drawing custom shapes
\usetikzlibrary{shadows}


%-----------------------------------------------------
%											 		\rem 
% insertion d'une remarque avec R
%---------------------------------------------------
\newcommand\rem[1]{%
   \marginpar{%
   \tikzpicture[baseline={(title.base)}]
      \node[inner sep=5pt,text width=4cm,drop shadow={shadow yshift=-5pt,shadow xshift=5pt,uMainColor},fill=white] (box) {\vskip5pt \nointerlineskip #1};
      \node[right=10pt,inner sep=0pt,outer sep=10pt] (title) at (box.north west) {\bfseries\color{uMainColor}Remarque};
      \draw[draw=uMainColor,very thick](title.west)--(box.north west)--(box.south west)--(box.south east)--(box.north east)--(title.east);
      \fill[uMainColor]([yshift=-10pt]box.north west)--+(-5pt,-5pt)--+(0pt,-10pt);
   \endtikzpicture}%
}
%-----------------------------------------------------
%											 \utikzimage 
% insertion d'une image sur fichier tkz.tex
%---------------------------------------------------

\newcommand{\utikzimage}[3]
{
\begin{wrapfigure}{R}{#3\textwidth}
  \begin{center}
	 \resizebox{#3\textwidth}{!}{\input{../#1.tkz.tex}}
  \end{center}
\caption{#2}
\end{wrapfigure}
}

%-----------------------------------------------------
%														\link 
% Url et footnote du lien
%-----------------------------------------------------

% TT every character and hyphenate after it
\def\hyphenateAndTtWholeString #1{\xHyphenate#1$\wholeString\unskip}

\def\xHyphenate#1#2\wholeString {\if#1$%
    \else\transform{#1}%
    \takeTheRest#2\ofTheString\fi}

\def\takeTheRest#1\ofTheString\fi
{\fi \xHyphenate#1\wholeString}

\def\transform#1{\url{#1}\hskip 0pt plus 1pt}

% Define the \urlx command which works like \url, but with line brakes
\def\urlx #1{\href{#1}{\hyphenateAndTtWholeString{#1}}}
  

 %-----------------------------------------------------
%														\link 
%-----------------------------------------------------  
  
  \newcommand{\ulink}[2]
  {\href{#1}
  {#2~\raisebox{-0.2ex}{\faExternalLink}\footnote{\url{#1}}}\xspace}
  
  \newcommand{\wikipedia}[1]
  { \href{https://fr.wikipedia.org/wiki/#1}{\textcolor{grey}{#1}}\raisebox{0.5ex}{\tiny{\textcolor{uMainColor}{\faWikipediaW}}}}
  


 %-----------------------------------------------------
%														\utodo 
%-----------------------------------------------------

  \newcommand{\utodo}{\begin{center}{\noindent \textcolor{cyan}{\faGears~{En cours de rédaction, DRAFT non publiable} \faGears}}\end{center}}
  
    \newcommand{\utocomplete}{\begin{center}{\noindent \textcolor{blue}{\faClipboard~{à compléter, les éléments ne donnent qu'une vue trop réduite ou trop parcellaire du sujet}}}\end{center}}
  

 




%-------------------------------------------------------------
%               FR CYBERDEF SECOPS COURSE
%
%                                    EDX STYLE
%
%                              2020 eduf@ction
%-------------------------------------------------------------
\usepackage{enumitem}

\usepackage{parskip}
\setlength{\parindent}{0em}
\setlength{\parskip}{5pt} % 1ex plus 0.5ex minus 0.2ex}

\def\uCoverColor{}

%---------------------------------------------------------------
%	MARGINS
%---------------------------------------------------------------

\usepackage{geometry}

\geometry{
	paper=a4paper, 
%	top=1cm, 
%	bottom=1cm, 
%	left=1cm, 
%	right=1cm, 
%	headheight=20pt, 
%	footskip=1.4cm, 
%	headsep=10pt, 
%	showframe, 
	showcrop 
}

%---------------------------------------------------------------
%	FOOTNOTE (format)
%---------------------------------------------------------------

\makeatletter
\long\def\@makefntextFB#1{%
    \ifx\thefootnote\ftnISsymbol
        \@makefntextORI{#1}%
    \else
        \rule\z@\footnotesep
        \setbox\@tempboxa\hbox{\@thefnmark}%
            \ifdim\wd\@tempboxa>\z@
                \kern2em\llap{\@thefnmark.\kern0.5em}%
            \fi
        \hangindent2em\hangafter\@ne#1
    \fi}
\makeatother

%-----------------------------------------------------
%	COLOR
%-----------------------------------------------------
%\usepackage{xcolor} % Required for specifying colors by name
%\definecolor{uMainColor}{RGB}{160,0,0} 
\colorlet{edxcolorcover}{uCoverColor}
%\definecolor{edxcolorcover}{RGB}{160,0,0} 

%\definecolor{cnam}{RGB}{128, 0, 32} 
                 

  \definecolor{comments}{rgb}{0.7,0,0}    % rouge foncé
  \definecolor{link}{rgb}{0,0.4,0.6}      % ~RoyalBlue de dvips
  \definecolor{url}{rgb}{0.6,0,0}         % rouge-brun
  \definecolor{citation}{rgb}{0,0.5,0}    % vert foncé
  \definecolor{ULlinkcolor}{rgb}{0,0,0.3} % de ulthese.cls
  \definecolor{rouge}{rgb}{0.85,0,0.07}   % rouge bandeau identitaire
  \definecolor{or}{rgb}{1,0.8,0}          % or bandeau identitaire

%-----------------------------------------------------
%	FONTS
%-----------------------------------------------------

%%\usepackage{avant} % Use the Avantgarde font for headings
%\usepackage{times} % Use the Times font for headings
%\usepackage{mathptmx} % Use the Adobe Times Rgman as the default text font together with math symbols from the Sym­bol, Chancery and Com­puter Modern fonts

%%\usepackage{microtype} % Slightly tweak font spacing for aesthetics
%%\usepackage[utf8]{inputenc} % Required for including letters with accents
%\usepackage[T1]{fontenc} % Use 8-bit encoding that has 256 glyphs

% Using AVANT GARDE Family

%\renewcommand{\familydefault}{\sfdefault}

%\fontfamily{pag}\selectfont
%\renewcommand{\familydefault}{pag} % sfdefault pag lmss

%bch         Charter
%lmr         Latin Modern Roman
%lmss        Latin Modern Sans Serif
%lmssq       Latin Modern Sans Serif extended
%lmtt        Latin Modern Typewriter
%lmvtt       Latin Modern Typewriter proportional
%pag         Avant Garde
%pbk         Bookman
%pcr         Courier
%phv         Helvetica
%pnc         New Century Schoolbook
%ppl         Palatino
%ptm         Times
%put         Utopia
 
%\renewcommand{\rmdefault}{pag} % text 
%\renewcommand{\sfdefault}{pag} % titre


%\renewcommand{\rmdefault}{pag} % text 
%\renewcommand{\sfdefault}{phv} % titre


%-----------------------------------------------------
%	BIBLIOGRAPHY AND INDEX
%-----------------------------------------------------


\usepackage{calc} % For simpler calculation - used for spacing the index letter headings correctly

\usepackage{imakeidx}
%\makeindex % Tells LaTeX to create the files required for indexing

\makeindex[columns=2,intoc=true,options={-s \upath/indexstyle.ist}]

%-----------------------------------------------------
%	HEADERS AND FOOTERS
%-----------------------------------------------------


\usepackage{fancyhdr} % Required for header and footer configuration
\pagestyle{fancy} % Enable the custom headers and footers

\ifcsname chapter \endcsname   % verify if command name exists in the CLASS (Book, report, article ...)
% BEGIN IF
\renewcommand{\chaptermark}[1]{\markboth{\sffamily\normalsize\bfseries\chaptername\ \thechapter.\ #1}{}} % Styling for the current chapter in the header
  \else
%nothing
\fi
% END IF

\renewcommand{\sectionmark}[1]{\markright{\sffamily\normalsize\thesection\hspace{5pt}#1}{}} % Styling for the current section in the header

\fancyhf{} % Clear default headers and footers
\fancyhead[LE,RO]{\sffamily\normalsize\thepage} % Styling for the page number in the header
\fancyhead[LO]{\sffamily\normalsize\rightmark} % Print the nearest section name on the left side of odd pages
\fancyhead[RE]{\sffamily\normalsize\leftmark} % Print the current chapter name on the right side of even pages
\fancyfoot[RE,LO]{\sffamily\normalsize\uCoursetittle} % Uncomment to include a footer
\fancyfoot[LE,RO]{\hyperlink{toc}{\sffamily\normalsize\uMinilogo}} % Uncomment to include a footer
% \hyperlink{toc}
\renewcommand{\headrulewidth}{1.0pt} % Thickness of the rule under the header

\fancypagestyle{plain}{% Style for when a plain pagestyle is specified
	\fancyhead{}\renewcommand{\headrulewidth}{0pt}} 

% Removes the header from odd empty pages at the end of chapters
\makeatletter
\renewcommand{\cleardoublepage}{
\clearpage\ifodd\c@page\else
\hbox{}
\vspace*{\fill}
\thispagestyle{empty}
\newpage
\fi}


%----------------------------------------------------------------------------------------
%	MAIN TABLE OF CONTENTS
%----------------------------------------------------------------------------------------
%\usepackage{titletoc} % Required for manipulating the table of contents

\contentsmargin{0cm} % Removes the default margin

% Part text styling (this is mostly taken care of in the PART HEADINGS section of this file)
\titlecontents{part}
	[0cm] % Left indentation
	{\addvspace{20pt}\bfseries} % Spacing and font options for parts
	{}
	{}
	{}

% Chapter text styling
\titlecontents{chapter}
	[1.25cm] % Left indentation
	{\addvspace{12pt}\large\sffamily\bfseries} % Spacing and font options for chapters
	{\color{uMainColor!60}\contentslabel[\Large\thecontentslabel]{1.25cm}\color{uMainColor}} % Formatting of numbered sections of this type
	{\color{uMainColor}} % Formatting of numberless sections of this type
	{\color{uMainColor!60}\normalsize\;\titlerule*[.5pc]{.}\;\thecontentspage} % Formatting of the filler to the right of the heading and the page number

% Section text styling
\titlecontents{section}
	[1.25cm] % Left indentation
	{\addvspace{3pt}\sffamily\bfseries} % Spacing and font options for sections
	{\contentslabel[\thecontentslabel]{1.25cm}} % Formatting of numbered sections of this type
	{} % Formatting of numberless sections of this type
	{\hfill\color{black}\thecontentspage} % Formatting of the filler to the right of the heading and the page number

% Subsection text styling
\titlecontents{subsection}
	[1.25cm] % Left indentation
	{\addvspace{1pt}\sffamily\small} % Spacing and font options for subsections
	{\contentslabel[\thecontentslabel]{1.25cm}} % Formatting of numbered sections of this type
	{} % Formatting of numberless sections of this type
	{\ \titlerule*[.5pc]{.}\;\thecontentspage} % Formatting of the filler to the right of the heading and the page number

%\titlecontents{subsubsection}
%	[1.25cm] % Left indentation
%	{\addvspace{1pt}\sffamily\tiny} % Spacing and font options for subsections
%	{\contentslabel[\thecontentslabel]{1.25cm}} % Formatting of numbered sections of this type
%	{} % Formatting of numberless sections of this type
%	{\ \titlerule*[.5pc]{.}\;\thecontentspage} % Formatting of the filler to the right of the heading and the page number

\titlecontents{subsubsection}
[1.25cm]
  {\addvspace{1pt}\sffamily\tiny}
 {\contentslabel[\thecontentslabel]{1.25cm}}
  {}
  {}

% Figure text styling
\titlecontents{figure}
	[1.25cm] % Left indentation
	{\addvspace{1pt}\sffamily\small} % Spacing and font options for figures
	{\thecontentslabel\hspace*{1em}} % Formatting of numbered sections of this type
	{} % Formatting of numberless sections of this type
	{\ \titlerule*[.5pc]{.}\;\thecontentspage} % Formatting of the filler to the right of the heading and the page number

% Table text styling
\titlecontents{table}
	[1.25cm] % Left indentation
	{\addvspace{1pt}\sffamily\small} % Spacing and font options for tables
	{\thecontentslabel\hspace*{1em}} % Formatting of numbered sections of this type
	{} % Formatting of numberless sections of this type
	{\ \titlerule*[.5pc]{.}\;\thecontentspage} % Formatting of the filler to the right of the heading and the page number

%----------------------------------------------------------------------------------------
%	MINI TABLE OF CONTENTS IN PART HEADS
%----------------------------------------------------------------------------------------

% Chapter text styling
\titlecontents{lchapter}
	[0em] % Left indentation
	{\addvspace{15pt}\large\sffamily\bfseries} % Spacing and font options for chapters
	{\color{uMainColor}\contentslabel[\Large\thecontentslabel]{1.25cm}\color{uMainColor}} % Chapter number
	{}  
	{\color{uMainColor}\normalsize\sffamily\bfseries\;\titlerule*[.5pc]{.}\;\thecontentspage} % Page number

% Section text styling
\titlecontents{lsection}
	[0em] % Left indentation
	{\sffamily\small} % Spacing and font options for sections
	{\contentslabel[\thecontentslabel]{1.25cm}} % Section number
	{}
	{}

% Subsection text styling (note these aren't shown by default, display them by searchings this file for tgcdepth and reading the commented text)
%\titlecontents{lsubsection}
%	[.5em] % Left indentation
%	{\sffamily\footnotesize} % Spacing and font options for subsections
%	{\contentslabel[\thecontentslabel]{1.25cm}}
%	{}
%	{}


\usepackage[]{ccicons}


%----------------------------------------------------------------------------------------
%	THEOREM STYLES
%----------------------------------------------------------------------------------------

\usepackage{amsmath,amsfonts,amssymb,amsthm} % For math equations, theorems, symbols, etc

\newcommand{\intoo}[2]{\mathopen{]}#1\,;#2\mathclose{[}}
\newcommand{\ud}{\mathop{\mathrm{{}d}}\mathopen{}}
\newcommand{\intff}[2]{\mathopen{[}#1\,;#2\mathclose{]}}
\renewcommand{\qedsymbol}{$\blacksquare$}

\ifcsname chapter \endcsname   % verify if command name exists in the CLASS (Book, report, article ...)
	\newtheorem{notation}{Notation}[chapter]
 \else
	\newtheorem{notation}{Notation}[section]
\fi


% Boxed/framed environments
\newtheoremstyle{ocrenumbox}% Theorem style name
{0pt}% Space above
{0pt}% Space below
{\normalfont}% Body font
{}% Indent amount
{\small\bf\sffamily\color{uMainColor}}% Theorem head font
{\;}% Punctuation after theorem head
{0.25em}% Space after theorem head
{\small\sffamily\color{uMainColor}\thmname{#1}\nobreakspace\thmnumber{\@ifnotempty{#1}{}\@upn{#2}}% Theorem text (e.g. Theorem 2.1)
\thmnote{\nobreakspace\the\thm@notefont\sffamily\bfseries\color{black}---\nobreakspace#3.}} % Optional theorem note

\newtheoremstyle{ocrenumbox}% Theorem style name
{0pt}% Space above
{0pt}% Space below
{\normalfont}% Body font
{}% Indent amount
{\small\bf\sffamily\color{uMainColor}}% Theorem head font
{\;}% Punctuation after theorem head
{0.25em}% Space after theorem head
{\small\sffamily\color{uMainColor}\thmname{#1}\nobreakspace\thmnumber{\@ifnotempty{#1}{}\@upn{#2}}% Theorem text (e.g. Theorem 2.1)
\thmnote{\nobreakspace\the\thm@notefont\sffamily\bfseries\color{black}---\nobreakspace#3.}} % Optional theorem note

\newtheoremstyle{blacknumex}% Theorem style name
{5pt}% Space above
{5pt}% Space below
{\normalfont}% Body font
{} % Indent amount
{\small\bf\sffamily}% Theorem head font
{\;}% Punctuation after theorem head
{0.25em}% Space after theorem head
{\small\sffamily{\tiny\ensuremath{\blacksquare}}\nobreakspace\thmname{#1}\nobreakspace\thmnumber{\@ifnotempty{#1}{}\@upn{#2}}% Theorem text (e.g. Theorem 2.1)
\thmnote{\nobreakspace\the\thm@notefont\sffamily\bfseries---\nobreakspace#3.}}% Optional theorem note


\newtheoremstyle{blacknumboxS}% name of the style to be used
{5pt}% measure of space to leave above the theorem. E.g.: 3pt
{5pt}% measure of space to leave below the theorem. E.g.: 3pt
{\normalfont}% name of font to use in the body of the theorem
{}% measure of space to indent
{\small\bf\sffamily}% name of head font
{}% punctuation between head and body
{ }% space after theorem head; " " = normal interword space
{\color{black}\faEye~  \color{uMainColor}\small\sffamily\thmnote{#3 : }}

% Non-boxed/non-framed environments
\newtheoremstyle{ocrenum}% Theorem style name
{5pt}% Space above
{5pt}% Space below
{\normalfont}% Body font
{}% Indent amount
{\small\bf\sffamily\color{uMainColor}}% Theorem head font
{\;}% Punctuation after theorem head
{0.25em}% Space after theorem head
{\small\sffamily\color{uMainColor}\thmname{#1}\nobreakspace\thmnumber{\@ifnotempty{#1}{}\@upn{#2}}% Theorem text (e.g. Theorem 2.1)
\thmnote{\nobreakspace\the\thm@notefont\sffamily\bfseries\color{black}---\nobreakspace#3.}} % Optional theorem note
\makeatother

% Defines the theorem text style for each type of theorem to one of the three styles above



\ifcsname chapter \endcsname

\newcounter{dummy} 

\numberwithin{dummy}{section}
	\theoremstyle{ocrenumbox}
		\newtheorem{theoremeT}[dummy]{Proposition}
		\newtheorem{problem}{Problem}[chapter]
		\newtheorem{exerciseT}{Outillage}[chapter]
	\theoremstyle{blacknumex}
		\newtheorem{exampleT}{Example}[chapter]
	\theoremstyle{blacknumbox}
		\newtheorem{vocabulary}{Vocabulary}[chapter]
		\newtheorem{definitionT}{Definition}[section]
		\newtheorem{corollaryT}[dummy]{Concept}
	\theoremstyle{blacknumboxS}
		\newtheorem{remarqueST}[]{}
		\newtheorem{remarqueT}[dummy]{Remarque}
	\theoremstyle{ocrenum}
	\newtheorem{proposition}[dummy]{Proposition}
  
 \else
 
\newcounter{dummy} 
\numberwithin{dummy}{section}

	\theoremstyle{ocrenumbox}
		\newtheorem{theoremeT}[dummy]{Proposition}
		%\newtheorem{problem}{Problem}[section]
		\newtheorem{exerciseT}{Outillage}[section]
	\theoremstyle{blacknumex}
		\newtheorem{exampleT}{Example}[section]
	\theoremstyle{blacknumbox}
		\newtheorem{vocabulary}{Vocabulary}[section]
		\newtheorem{definitionT}{Definition}[section]
		\newtheorem{corollaryT}[dummy]{concept}
	\theoremstyle{blacknumboxS}
		\newtheorem{remarqueST}[]{}
		\newtheorem{remarqueT}[dummy]{Remarque}
	\theoremstyle{ocrenum}
		\newtheorem{proposition}[dummy]{Proposition}
\fi



%----------------------------------------------------------------------------------------
%	DEFINITION OF COLORED BOXES
%----------------------------------------------------------------------------------------

\RequirePackage[framemethod=default]{mdframed} % Required for creating the theorem, definition, exercise and corollary boxes

% Theorem box
\newmdenv[skipabove=7pt,
skipbelow=7pt,
backgroundcolor=black!5,
linecolor=uMainColor,
innerleftmargin=5pt,
innerrightmargin=5pt,
innertopmargin=5pt,
leftmargin=0cm,
rightmargin=0cm,
innerbottommargin=5pt]{tBox}

% Exercise box	  
\newmdenv[skipabove=7pt,
skipbelow=7pt,
rightline=false,
leftline=true,
topline=false,
bottomline=false,
backgroundcolor=uMainColor!10,
linecolor=uMainColor,
innerleftmargin=5pt,
innerrightmargin=5pt,
innertopmargin=5pt,
innerbottommargin=5pt,
leftmargin=0cm,
rightmargin=0cm,
linewidth=4pt]{eBox}	

% Definition box
\newmdenv[skipabove=7pt,
skipbelow=7pt,
rightline=false,
leftline=true,
topline=false,
bottomline=false,
linecolor=uMainColor,
innerleftmargin=5pt,
innerrightmargin=5pt,
innertopmargin=0pt,
leftmargin=0cm,
rightmargin=0cm,
linewidth=4pt,
innerbottommargin=0pt]{dBox}	

% Corollary box
\newmdenv[skipabove=7pt,
skipbelow=7pt,
rightline=false,
leftline=true,
topline=false,
bottomline=false,
linecolor=gray,
backgroundcolor=black!5,
innerleftmargin=5pt,
innerrightmargin=5pt,
innertopmargin=5pt,
leftmargin=0cm,
rightmargin=0cm,
linewidth=4pt,
innerbottommargin=5pt]{cBox}

% Creates an environment for each type of theorem and assigns it a theorem text style from the "Theorem Styles" section above and a colored box from above
%\newenvironment{theorem}{\begin{tBox}\begin{theoremeT}}{\end{theoremeT}\end{tBox}}
\newenvironment{exercise}{\begin{eBox}\begin{exerciseT}}{\hfill{\color{uMainColor}\tiny\ensuremath{\blacksquare}}\end{exerciseT}\end{eBox}}				  
%\newenvironment{definition}{\begin{dBox}\begin{definitionT}}{\end{definitionT}\end{dBox}}	
%\newenvironment{example}{\begin{exampleT}}{\hfill{\tiny\ensuremath{\blacksquare}}\end{exampleT}}		
%\newenvironment{corollary}{\begin{cBox}\begin{corollaryT}}{\end{corollaryT}\end{cBox}}	
\newenvironment{nota}{\begin{cBox}\begin{remarqueT}}{\end{remarqueT}\end{cBox}}	

%----------------------------------------------------------------------------------------
%	REMARK ENVIRONMENT
%----------------------------------------------------------------------------------------

\newenvironment{remark}{\par\vspace{10pt}\small % Vertical white space above the remark and smaller font size
\begin{list}{}{
\leftmargin=35pt % Indentation on the left
\rightmargin=25pt}\item\ignorespaces % Indentation on the right
\makebox[-2.5pt]{\begin{tikzpicture}[overlay]
\node[draw=uMainColor!60,line width=1pt,circle,fill=uMainColor!25,font=\sffamily\bfseries,inner sep=2pt,outer sep=0pt] at (-15pt,0pt){\textcolor{uMainColor}{i}};\end{tikzpicture}} % Orange R in a circle
\advance\baselineskip -1pt}{\end{list}\vskip5pt} % Tighter line spacing and white space after remark

%----------------------------------------------------------------------------------------
%	SECTION NUMBERING IN THE MARGIN
%----------------------------------------------------------------------------------------

\makeatletter
% \renewcommand{\@seccntformat}[1]{\textcolor{uMainColor}{\csname the#1\endcsname}\quad}   

%\makeatother

%----------------------------------------------------------------------------------------
%	PART HEADINGS
%----------------------------------------------------------------------------------------

% Numbered part in the table of contents
\newcommand{\@mypartnumtocformat}[2]{%
	\setlength\fboxsep{0pt}%
	\noindent\colorbox{uMainColor!20}{\strut\parbox[c][.7cm]{\ecart}{\color{uMainColor!70}\Large\sffamily\bfseries\centering#1}}\hskip\esp\colorbox{uMainColor!40}{\strut\parbox[c][.7cm]{\linewidth-\ecart-\esp}{\Large\sffamily\centering#2}}%
}

% Unnumbered part in the table of contents
\newcommand{\@myparttocformat}[1]{%
	\setlength\fboxsep{0pt}%
	\noindent\colorbox{uMainColor!40}{\strut\parbox[c][.7cm]{\linewidth}{\Large\sffamily\centering#1}}%
}

\newlength\esp
\setlength\esp{4pt}
\newlength\ecart
\setlength\ecart{1.2cm-\esp}
\newcommand{\thepartimage}{}%
\newcommand{\partimage}[1]{\renewcommand{\thepartimage}{#1}}%
\def\@part[#1]#2{%
\ifnum \c@secnumdepth >-2\relax%
\refstepcounter{part}%
\addcontentsline{toc}{part}{\texorpdfstring{\protect\@mypartnumtocformat{\thepart}{#1}}{\partname~\thepart\ ---\ #1}}
\else%
\addcontentsline{toc}{part}{\texorpdfstring{\protect\@myparttocformat{#1}}{#1}}%
\fi%
\startcontents%
\markboth{}{}%
{\thispagestyle{empty}%
\begin{tikzpicture}[remember picture,overlay]%
\node at (current page.north west){\begin{tikzpicture}[remember picture,overlay]%	
\fill[uMainColor!20](0cm,0cm) rectangle (\paperwidth,-\paperheight);
\node[anchor=north] at (4cm,-3.25cm){\color{uMainColor!40}\fontsize{220}{100}\sffamily\bfseries\thepart}; 
\node[anchor=south east] at (\paperwidth-1cm,-\paperheight+1cm){\parbox[t][][t]{8.5cm}{
\printcontents{l}{0}{\setcounter{tocdepth}{1}}% The depth to which the Part mini table of contents displays headings; 0 for chapters only, 1 for chapters and sections and 2 for chapters, sections and subsections
}};
\node[anchor=north east] at (\paperwidth-1.5cm,-3.25cm){\parbox[t][][t]{15cm}{\strut\raggedleft\color{white}\fontsize{30}{30}\sffamily\bfseries#2}};
\end{tikzpicture}};
\end{tikzpicture}}%
\@endpart}
\def\@spart#1{%
\startcontents%
\phantomsection
{\thispagestyle{empty}%
\begin{tikzpicture}[remember picture,overlay]%
\node at (current page.north west){\begin{tikzpicture}[remember picture,overlay]%	
\fill[uMainColor!20](0cm,0cm) rectangle (\paperwidth,-\paperheight);
\node[anchor=north east] at (\paperwidth-1.5cm,-3.25cm){\parbox[t][][t]{15cm}{\strut\raggedleft\color{white}\fontsize{30}{30}\sffamily\bfseries#1}};
\end{tikzpicture}};
\end{tikzpicture}}
\addcontentsline{toc}{part}{\texorpdfstring{%
\setlength\fboxsep{0pt}%
\noindent\protect\colorbox{uMainColor!40}{\strut\protect\parbox[c][.7cm]{\linewidth}{\Large\sffamily\protect\centering #1\quad\mbox{}}}}{#1}}%
\@endpart}
\def\@endpart{\vfil\newpage
\if@twoside
\if@openright
\null
\thispagestyle{empty}%
\newpage
\fi
\fi
\if@tempswa
\twocolumn
\fi}

%----------------------------------------------------------------------------------------
%	CHAPTER HEADINGS
%-------------------------------------------------%--------------------------------------

%
%
%% A switch to conditionally include a picture, implemented by Christian Hupfer
\newif\ifusechapterimage
\usechapterimagetrue
\newcommand{\thechapterimage}{}%
\newcommand{\chapterimage}[1]{\ifusechapterimage\renewcommand{\thechapterimage}{#1}\fi}%
\newcommand{\autodot}{.}

\def\@makechapterhead#1{%
{\parindent \z@ \raggedright \normalfont
\ifnum \c@secnumdepth >\m@ne
\if@mainmatter
\begin{tikzpicture}[remember picture,overlay]
	
\node at (current page.north west)
			{\begin{tikzpicture}[remember picture,overlay]
			\node[anchor=north west,inner sep=0pt] at (0,0) {\ifusechapterimage\includegraphics[width=\paperwidth]{\thechapterimage}\fi};
			\draw[anchor=west] (\Gm@lmargin,-9cm) node [line width=2pt,rounded corners=15pt,draw=uMainColor,fill=white,fill opacity=1,inner sep=15pt]{\strut\makebox[22cm]{}};
			
			\draw[anchor=west] (\Gm@lmargin+.3cm,-9cm) node {\fontsize{20}{30}\selectfont\sffamily\bfseries\color{uMainColor}~#1\strut};
			%{\Huge\sffamily\bfseries\color{uMainColor}\thechapter\autodot~#1\strut};
			\node at (current page.north east)  [xshift=-3.95cm, yshift=-4.05cm, text opacity=1]  {{\color{white}\centering\fontsize{200}{30}\selectfont \bfseries\sffamily\thechapter\strut} };
			\node at (current page.north east)  [xshift=-4cm, yshift=-4cm, text opacity=1]  {{\color{uMainColor}\centering\fontsize{200}{30}\selectfont \bfseries\sffamily\thechapter\strut} };
			
			\end{tikzpicture}};
	\end{tikzpicture}

\else

\begin{tikzpicture}[remember picture,overlay]
\node at (current page.north west)
{\begin{tikzpicture}[remember picture,overlay]
\node[anchor=north west,inner sep=0pt] at (0,0) {\ifusechapterimage\includegraphics[width=\paperwidth]{\thechapterimage}\fi};
\draw[anchor=west] (\Gm@lmargin,-9cm) node [line width=2pt,rounded corners=15pt,draw=uMainColor,fill=white,fill opacity=1,inner sep=15pt]{\strut\makebox[22cm]{}};

\draw[anchor=east] (\Gm@lmargin,-9cm) node [line width=2pt,rounded corners=15pt,draw=uMainColor,fill=white,fill opacity=1,inner sep=15pt]{\strut\makebox[22cm]{}};
\draw[anchor=west] (\Gm@lmargin+.3cm,-9cm) node {\huge\sffamily\bfseries\color{uMainColor}#1\strut};

\end{tikzpicture}};
 
\end{tikzpicture}
\fi\fi\par\vspace*{270\p@}}}

\makeatother
%-------------------------------------------

%\def\@makeschapterhead#1{%
%\begin{tikzpicture}[remember picture,overlay]
%\node at (current page.north west)
%{\begin{tikzpicture}[remember picture,overlay]
%\node[anchor=north west,inner sep=0pt] at (0,0) {\ifusechapterimage\includegraphics[width=\paperwidth]{\thechapterimage}\fi};
%\draw[anchor=west] (\Gm@lmargin,-9cm) node [line width=2pt,rounded corners=15pt,draw=uMainColor,fill=white, fill opacity=0.8,inner sep=15pt]{\strut\makebox[22cm]{}};
%\draw[anchor=west] (\Gm@lmargin+.3cm,-9cm) node {\huge\sffamily\bfseries\color{uMainColor}#1\strut};
%\end{tikzpicture}};
%\end{tikzpicture}
%\par\vspace*{270\p@}}




%----------------------------------------------------------------------------------------
%	LINKS
%----------------------------------------------------------------------------------------


\usepackage{bookmark}
\bookmarksetup{
open,
numbered,
addtohook={%
\ifnum\bookmarkget{level}=0 % chapter
\bookmarksetup{bold}%
\fi
\ifnum\bookmarkget{level}=-1 % part
\bookmarksetup{color=uMainColor,bold}%
\fi
}
}

%----------------------------------------------------------------------------------------
%	FANCY CHAPTER Header
%----------------------------------------------------------------------------------------





%\definecolor{uMainColor}{RGB}{255,128,0}   



%-------------------------------------------------
%	BIBLIOGRAPHY
%-------------------------------------------------

\addbibresource{../Latex/Sources/commons/bibliography.bib} % BibTeX bibliography file
\defbibheading{bibempty}{}

%-------------------------------------------------
%	COLORS & BORDERS
%-------------------------------------------------
\setlength{\fboxrule}{0.75pt} % Width of the border around the abstract
%\definecolor{uMainColor}{RGB}{255,128,0} % Color of the article title and section
\definecolor{color2}{RGB}{220,220,220} % Color of the boxes behind the abstract and headings

%-------------------------------------------------
% ITEMS DEFINITION
%-------------------------------------------------

\setlist [itemize,1]{label=\color{uMainColor}\faCaretRight }


\usepackage{environ}
\newwrite\audioexport



\begin{document}





  %-------------------------------------------------------
  % AudioNote Open File
  %-------------------------------------------------------
\mode<all>{\immediate\openout\audioexport=\jobname.txt}


\addtocontents{toc}{\protect\hypertarget{toc}{}}

%-------------------------------------------------
%	ARTICLE INFORMATION
%-------------------------------------------------

\Abstract {\uabstract%===========================
% COURS "INTRO CYBERDEF"
% Abstract général des articles
%===========================


Il fait partie du cours introductif aux fondamentaux de la sécurité des systèmes d'information vue sous deux prismes quelques fois opposés dans la littérature : la gouvernance et la gestion opérationnelle de la sécurité.
Le cours est constitué d'un ensemble de notes de synthèse indépendantes compilées en un document unique, mais édité par chapitre dans le cadre de ce cours.\\
C\edoc ne constitue pas à lui seul le référentiel du cours CYBERDEF101 (SEC101 du Cnam). Il compile des notes de cours mises à disposition de l'auditeur comme support pédagogique partiel à ce cours introductif à la cyberdéfense d'entreprise. }
\JournalInfo{\uJournalInfo} % Journal information
\Archive{Notes de cours éditées le  \DTMnow} % Additional notes (e.g. copyright, DOI, review/research article)
\PaperTitle{\utitle} % Article title
\Authors{\uauthor\textsuperscript{1,}\textsuperscript{2}*} % Authors
\affiliation{\textsuperscript{1}\textit{\uproa}} % Author affiliation
\affiliation{\textsuperscript{2}\textit{\uprob}} % Author affiliation
\affiliation{*\textbf{email}: \umaila\ -- \umailb} % Corresponding author
\Keywords{\ukeywords} % Keywords 
\newcommand{\keywordname}{Mots clefs} % Defines the keywords heading name

%-------------------------------------------------
%	MAKE TITLE
%-------------------------------------------------

\maketitle

\begin{tikzpicture}[remember picture,overlay]\node at (current page.north west)[xshift=3cm, yshift=-3cm, text opacity=1.0]{\includegraphics[width=0.1\paperwidth]{\utemplatepath/commons.inc/CommonsPictures/\utemplatetheme.theme/shield-20.pdf}}; 
\end{tikzpicture} 


\AddToShipoutPicture{%
\AtPageCenter{%
\begin{tikzpicture}[remember picture,overlay]
\node [rotate=60,scale=5,text opacity=0.1] at (current page.center) { \textit{\color{uAlertColor}\INFODistrib} };
\end{tikzpicture} 
%\begin{tikzpicture}[remember picture,overlay]
%\node [rotate=60,scale=10,text opacity=0.3]
%at (current page.center) {moi};
%\end{tikzpicture}
}%
}


%-------------------------------------------------	
\begin{center}
	\textcolor{uMainColor}{\textbf{\INFODistrib}} % Distribution information
\end{center}

\INFOGitHub % Link to download file and copyright
%-------------------------------------------------

%-------------------------------------------------
	\newpage
	\tableofcontents  % TOC
	\iftotalfigures
  			\listoffigures
	\fi
%-------------------------------------------------

%-------------------------------------------------
\newpage
\ubody % MAIN BODY defined external
%-------------------------------------------------

	\printbibliography
 \printglossary[type=\acronymtype]
%	\printglossary
%-------------------------------------------------	
\newpage
\Ucontribute % Contribution TEXT
%-------------------------------------------------

  %-------------------------------------------------------
  % AudioNote Open File
  %-------------------------------------------------------
  \mode<all>{\immediate\closeout\audioexport}

\end{document}
}
	\fi
\fi


 
\lessonID{7.0}
%************************************************************
% Chargement des variables dédiées à l'article
%************************************************************

\setauthor{EDU}

%========================================
\newcommand {\ukeywords}
%========================================
{%---------------------------------------------------------------------
Anticipation, veille, alerte, réponse, CERT, SOC
}%

%========================================
\newcommand {\utitle}
{%---------------------------------------------------------------------
Introduction à la SECOPS
}%---------------------------------------------------------------------

%========================================
\newcommand {\uabstract}
{%---------------------------------------------------------------------
C\edoc introduit le triptype de la partie cyberdéfense de la sécurité opérationnelle : Anticiper, Détecter, Réagir  et ceci sur les trois grands invariants des risques numériques : les vulnérabilités, les menaces et l'impact.
}
%---------------------------------------------------------------------

%************************************************************
%  variable définissant  le corps de l'article
%************************************************************

%========================================
\newcommand {\ubody}
{%---------------------------------------------------------------------
%$BEGIN_UBODY


\mode<all>{%-------------------------------------------------------------
%               FR CYBERDEF SECOPS COURSE
%                                        SECOPS
%.                    Vulnerability/Threat/Incident
%
%                                    Introduction
%                                chap-VTI-intro.tex
%
%                              2020 eduf@ction
%-------------------------------------------------------------
\uchap{Chapitre introductif de la Partie 3 du cours SEC 101 (VTI-INTRO)}


%=======================================
\section{Sécurité opérationnelle}
%--------------------------------------------------------------

En terme de gouvernance, après avoir construit une structure de sécurité cohérente sur les aspects de gestion des flux, de gestion des accès et des identités, et construit une gouvernance efficace sur la base de l'\uindex{ISO27001}, il est  nécessaire de maintenir le niveau de sécurité de l'entreprise ou du système. La dynamique de sécurité de l'entreprise en exploitation nécessite une organisation et des politiques orientées vers cette efficience de l'anticipation, de la détection et de la réaction que nous avons présentées comme le volet cyberdéfense de la cybersécurité de l'entreprise.


% begin frame 
\mode<all>{\texframe{sécurité opérationnelle}{Dans le cycle de vie}
{
Dans certains ouvrages ce processus est dénommé  \g{Maintien en Condition de sécurité}. En utilisant  les termes anglo-saxons définissant le cycle des projets, nous pourrions positionner ses activités dans la phase dite de  \textbf{RUN}. Les autres phases en amont pouvant être définies comme : 

\begin{itemize}
 \item \textbf{THINK/DESIGN} : Des risques évalués à la politique sécurité établie en fonction des risques ;
 \item \textbf{BUILD}: De la politique de sécurité déployée à la construction d’une sécurité implémentée;
\end{itemize}
}} % end frame




% begin PRZ==========================
\begin{frame}
\frametitle<presentation>{sécurité opérationnelle}
\framesubtitle<presentation>{l'activité RUN}
% end header PRZ======================
\upicture{Latex/Sources/EDU/Pictures/img-thinkrun}{les phases du cycle de vie}{0.7}{lblthinkrun}

Et nous classons donc dans la dernière phase du cycle de vie : les activités d'exploitation de la sécurité, \textbf{RUN} : Des évènements de sécurité gérés à une \textbf{Cybercrise} maitrisée, ce que certains appellent \textbf{\uindex{SECOPS}} : \g{\uindex{sécurité Opérationnelle}}.
\end{frame}
% end PRZ ===========================


% begin PRZ==========================
\begin{frame}
\frametitle<presentation>{Objectifs adaptés aux finalités}
\framesubtitle<presentation>{de l'activité de l'entreprise}
% end header PRZ======================
Ce modèle se développe bien entendu en fonction des finalités de l'entreprise.

\begin{itemize}
    \item Soit nous sommes dans \textbf{une dynamique entreprise} et ces processus sont ceux mis en place pour s'assurer que  l'ensemble de actions sont prises en compte pour maîtriser les fragilités dont les vulnérabilités informatiques, détecter les menaces tant en anticipation que pendant des attaques (bruyantes, ou discrètes), et réagir pour maintenir l'activité et limiter l'impact.
    \item  Soit nous sommes \textbf{fabricant d'un produit ou d'un service}, et au delà des engagements sécuritaires de toute entreprise (cf. ci-dessus) des processus de \g{maintien en condition de sécurité} des produits et services sont à ajouter pour maîtriser les vulnérabilités, les correctifs et leur cycle de vie  (audit, communication, gestion des découvertes de fragilités par des tiers, rémunération de BugHunters ...).
\end{itemize}
\end{frame}
% end PRZ ===========================

% begin PRZ==========================
\begin{frame}
\frametitle<presentation>{sécurité opérationnelle}
\framesubtitle<presentation>{Une définition}
% end header PRZ======================
 Le terme de \g{sécurité opérationnelle}, est relativement jeune dans l’histoire de la sécurité des technologies de l’information. Le terme de \uindex{SSI}(sécurité des Systèmes d’Information) était né pour distinguer des disciplines qui s’attachaient à protéger l’information qui circulent dans les systèmes d'information de l’entreprise (cf. protection et classification de l’information) vis à vis de la sécurité des biens et des personnes. La sécurité des réseaux et la sécurité informatique ont été les précurseurs de la cybersécurité, le cyber recouvrant en un seul terme, les enjeux de sécurité liés au réseau et à l’informatique, mais plus largement à la sécurité de l'économie numérique.
\end{frame}
% end PRZ ===========================

Comme nous l’avons abordé dans l'introduction et dans les chapitres précédents, la cybersécurité est un domaine vaste qui regroupe de nombreuses disciplines. Elle peuvent intervenir dans des cycles projets pour construire des systèmes sûrs ou pour assurer la continuité d’activité et la protection des patrimoines dans l'entreprise. Il faut aussi, penser à y ajouter aussi tout un espace de gestion des conformités (législatives, réglementaires, normatives, contractuelles).


C’est plutôt dans un contexte d'opération sécurité au quotidien que l’on parle de sécurité opérationnelle. Ces activités opérationnelles supportent donc le maintien en condition de sécurité au quotidien de l’entreprise. En France, au sein des armées, on parle de lutte informatique défensive permettant de différencier les activités des Cyber-défense des activités de Cyber-protection. Ces activités sont à opposer à la lutte informatique\g {offensive} qui ne sera pas abordé dans c\ecours car elle relève de prérogative des états et non des entreprises. Nous aurons toutefois l'occasion d'aborder le \textbf{Hackback}, dans un prochain.
La sécurité opérationnelle ajoute par ailleurs à son périmètre de surveillance de l'intérieur de la zone de responsabilité de l'entreprise SI, réseaux sociaux, services cloud ...), celle extérieure à l'entreprise via des mécanismes de veille sur la menace et de surveillance des compromissions potentielles. Nous pourrions évoqué l'image d'une cité où \g{Les murs sont épais et solides, les douaniers sont aux portes de la cité, la police doit toutefois veiller à la sécurité des biens et des citoyens dans la ville, car certains sont néanmoins des brigands. Quand à l'armée, elle veille aux frontières du pays informée par des agents à l'étranger}.

On traitera donc cette partie avec une équivalence dans les terminologies suivantes :

% begin PRZ==========================
\begin{frame}
\frametitle<presentation>{Plusieurs terminologies, une dynamique}
% end header PRZ======================
	\begin{itemize}
		\item Maintien en condition de sécurité (\uindex{MCS}); 
		\item Sécurité opérationnelle (\uindex{SECOPS});
		\item Lutte informatique défensive (\uindex{LID});
		\item Cyberdéfense au sens de la cyberdéfense d'entreprise (CYBERDEFENSE).
	\end{itemize}
\end{frame}
% end PRZ ===========================

Le but de cette sécurité opérationnelle est d’être au coeur de l’action de la sécurité de l’entreprise. En effet, la sécurité de l’entreprise est une propriété multiforme. Elle est d’abord statique dans la mesure où elle correspond à un niveau de confiance dans l’environnement pour conserver la disponibilité, la confidentialité et l’intégrité de l’entreprise. Cette forme statique est souvent liée à la conformité de l’entreprise, aux différents référentiels sécuritaires (ISO27000, GDPR, LPM, NIS ...), mais surtout aux objectifs sécurité de l'entreprise face à ses risques et aux exigences de sécurité des clients souvent inscrites dans des plans d’assurance sécurité (Cf. PAS). 
Elle est aussi dynamique car c’est aussi une propriété systémique qui mesure la capacité à anticiper les menaces, identifier les fragilités , détecter en temps réel les attaques et réagir à temps ou au pire disposer des capacités de revenir dans un état de fonctionnement compatible avec la survie de l’entreprise (Modes dégradés par exemple).
Le système évolue, faisant apparaitre ici et là de nouvelles fragilités, l’entreprise se transforme, vit, suscitant de nouveaux potentiels d’attaques. L’entreprise doit s’organiser pour disposer de fonctions opérationnelles adaptées et dédiées à cette activité. Ces fonctions nécessitent des savoirs, des savoirs-faire et de l’outillage. C’est l’ensemble de ces techniques que nous allons tenter d’aborder dans c\edoc.

Globalement, on peut remarquer que le cycle de vie est à prendre dans le sens inverse de notre présentation. Dans les entreprises moins matures en gouvernance de la sécurité, la dynamique de cette sécurité opérationnelle est la première visible et opérée. L'entreprise va réagir le plus souvent dans une dynamique de réponse immédiates aux problèmes de sécurité sans pour autant investiguer plus avant dans les fragilités globales. Les mécanismes de cybersécurité sont donc construits dans une entreprise peu mature dans le sens suivant :

% begin PRZ==========================
\begin{frame}
\frametitle<presentation>{Enjeux SECOPS}
% end header PRZ======================
	\begin{itemize}
		\item Répondre aux incidents de sécurité, tenter de répondre à la question : \g{qui nous attaque et pourquoi}; 
		\item Améliorer les filtrages;
		\item Couvrir les vulnérabilités découvertes;
		\item Rechercher les vulnérabilités existantes dans le périmètre de responsabilité ;
		\item Anticiper les attaques;
		\item Anticiper les risques informatiques;
		\item Anticiper les risques sur l'information;
		\item Anticiper la menace.
	\end{itemize}
\end{frame}
% end PRZ ===========================

\section{Lutte contre la menace}

La finalité de cette défense d’entreprise est de lutter contre ces attaques qui ne sont pas qu’informatiques. L’attaquant peut utiliser des scenarii utilisant de nombreux vecteurs qui peuvent utiliser des fragilités organisationnelles ou humaines. On peut dire qu’une attaque est une fonction complexe, qui peut viser ou utiliser de nombreux facteurs internes et externes à l’entreprise. Ces facteurs constituent ce que certains nomment l’environnement numérique ou digital de l’entreprise. Cet environnement est globalement constitué de l’ensemble des outils, services, moyens informatiques ou réseaux utilisés par l’entreprise.
Mai 2017 a été un tournant dans la prise de conscience de la menace de la part des entreprises. Le Rançon-logiciel WannaCry a plus fait trembler les médias que les entreprises, mais a permis de faire comprendre au grand public les enjeux des menaces informatiques.


\begin{nota}[Paramètres d’une attaque]
\begin{equation}
Attaque = Fonction \left[ Fragilit\acute{e}s\,HOT\, Entreprise\otimes Gains\,Escompt\acute{e}s\,PF \right]
\end{equation}
\end{nota}

\begin{itemize}
	\item Fragilités HOT : Humaines, Organisationnelles, Techniques 
	\item Gains pour l'attaquant : Idéologiques, Politiques, Financiers, ...
\end{itemize}

On peut classer la majorité des attaques informatiques dans quatre grandes classes :

% begin PRZ==========================
\begin{frame}
\frametitle<presentation>{Grandes typologies des attaques numériques}
% end header PRZ======================
\begin{itemize}
\item Attaques \textbf{d’interception} d’information, vols par écoutes passives ou actives dans les flux transitant entre un émetteur et un récepteur;
\item Attaques par \textbf{déni de services}, généralement sur le réseau : Ce type d’attaque est une atteinte à la DISPONIBILITE du système, basé souvent sur la saturation d’une capacité de traitement. Le système saturé dans l’exécution de certaines de ses fonctions, ne peut plus répondre aux demandes légitimes, car il est occupé à traiter d’autres sollicitations;
\item Attaques par \textbf{exploitation de failles }logiciels : Ce type d’attaque va utiliser une vulnérabilité, d’un système d’exploitation ou d’un logiciel pour exécuter du code malveillant. Ce code réalisera alors sa mission;
\item Attaques par \textbf{exploitation de défauts} de configuration : Ce type d’attaque utilise simplement un ou des défauts de configuration pour que légitimement l’agresseur puisse dérouler un scénario, qui pourra lui donner par exemple des droits particuliers pour conduire des attaques.
\end{itemize}
\end{frame}
% end PRZ==========================

Nous pourrions remarquer que ce nombre est relativement faible. Toutefois, la vrai difficulté réside dans la multiplicité des vulnérabilités, et des défauts de configuration. Les développeurs réalisent des logiciels possédant des failles (vulnérabilités), les utilisateurs ou les administrateurs déploient des systèmes en faisant des erreurs de configuration, ou ne les configurent que très rarement en pensant à la malveillance.

Les motivations des attaquants sont nombreuses, et leurs objectifs variés :

% begin PRZ==========================
\begin{frame}
\frametitle<presentation>{Motivations de l'attaquant}
% end header PRZ======================

\begin{itemize}
\item obtenir un accès au système pour s’y maintenir en attendant une opportunité ;
\item récupérer de l’information, secrets, données personnelles exploitables (en gros toute information ayant de la valeur)
\item récupérer des données bancaires ;
\item s'informer sur l'organisation (entreprise de l'utilisateur, etc.) ;
\item troubler, couper, bloquer le fonctionnement d'un service (les \uindex{rançongiciels} entre dans cette catégories) ;
\item utiliser le système d’un utilisateur, pour rebondir vers un autre système ;
\item détourner les ressources du système d’un utilisateur (utiliser de la bande passante, utiliser de la capacité de calcul) ;
\end{itemize}
\end{frame}
% end PRZ==========================

Bien entendu, il n’y a que très rarement un seul objectif, c’est la combinaison des méthodes d’attaques, des objectifs unitaires qui définissent globalement une mission ou un objectif final. L’exploitation de vulnérabilités au sein de l’entreprise va permettre le déploiement par l’attaquant d’un scénario.

\subsection{Politiques et Stratégies}

% begin PRZ==========================
\begin{frame}
\frametitle<presentation>{Politique vs stratégie}
% end header PRZ======================

\upicture{Latex/Sources/EDU/Pictures/img-cyclevie-pol-strat}{Positionnement de la sécurité opérationnelle}{0.9}{lblpol-strat}

\end{frame}
% end PRZ==========================

A partir des risques identifiés, l’entreprise a posé des politiques de sécurité qui ont permis de mettre en place des mesures de sécurité. Ces mesures sont d’ordre techniques avec des systèmes de sécurité, ou des SI avec des architectures particulières, mais aussi d’ordre organisationnel avec des procédures et des mécanismes à respecter.
L’ensemble de cette dynamique construit un niveau de sécurité qu’il va être nécessaire de maintenir dans le temps. Toutefois ce niveau de sécurité n’est pas suffisant pour une simple et bonne raison : la menace évolue, les vulnérabilités apparaissent (découvertes, ou créées), la valeur\g{marchande} des actifs d’une entreprise change aussi. 
Les occurrences de ces éléments de vie sont considérés comme des évènements qu’il convient de détecter avec suffisamment d’avance sur l’attaquant pour pouvoir le plus rapidement les prendre en compte.

La gestion des événements qui peuvent être une source de mesure de l’évolution du niveau de sécurité de l’entreprise est au coeur des stratégies de cyberdéfense. Ces évènements sont corrélés avec des sources provenant de deux processus particuliers qui seront décrits dans ce c\edoc.

% begin PRZ==========================
\begin{frame}
\frametitle<presentation>{Politique vs stratégie}
% end header PRZ======================
Il est à noter qu'un attaquant ne raisonne pas en politiques d'attaque face à une politique de sécurité, mais par des stratégies auxquelles il faut opposer aussi par des stratégies de défense, dont 

\begin{itemize}
\item Recherche des vulnérabilités : Processus qui permet de rechercher, découvrir, couvrir les vulnérabilités ou fragilités de l’entreprise ou ayant un impact sur l’entreprise que celles-ci soient techniques, humaines ou organisationnelles ;
\item Prévention de la menace : Processus qui permet de connaître les menaces directes sur l’entreprise ou potentielles afin d’anticiper et/ou se préparer à un type d’attaque.
\end{itemize}
\end{frame}
% end PRZ==========================

C’est la confrontation entre les vulnérabilités, les menaces et la détection de l’activité de l’entreprise qui va permettre d’être efficace dans le processus de réponse. Il y a de nombreuses manière d’aborder la cyberdéfense d’entreprise.

C\edoc présente donc une dynamique de cyberdéfense en trois\g{volets }

% begin PRZ==========================
\begin{frame}
\frametitle<presentation>{3 volets d'une cyberdéfense}
% end header PRZ======================


\begin{itemize}
\item Gestion des vulnérabilités (\textit{Vulnerability Management and CERT}) : maîtriser ses vulnérabilités mais aussi surveiller l’environnement technologique. 
\item Surveillance, Détection de la menace (\textit{Event and Threat Management}) : Analyser en temps réel l’environnement protégé mais aussi surveiller l’écosystème lié à la menace pour anticiper 
\item Gestion des incidents et réponse aux incidents (\textit{Incident Response – CSIRT}) : Réagir en cas d’incident et assurer la remédiation
\end{itemize}

\end{frame}
% end PRZ==========================


% begin PRZ==========================
\begin{frame}
\frametitle<presentation>{3 volets d'une cyberdéfense}
% end header PRZ======================
\upicture{Latex/Sources/EDU/Pictures/img-triangle}{Des 3 des volets de la sécurité opérationnelle}{0.6}{lbltriptyque}

\end{frame}
% end PRZ==========================

Ces trois volets ne sont pas les seuls qui concourent à la cyberdéfense d’entreprise, mais il en reste les trois faces principales. Il est à noter que ces trois volets correspondent aussi en France à trois référentiels de qualification de l’ANSSI des prestataires de services de cybersécurité au profit des entreprises. Ces labels sont obtenus par les entreprises qui respectent un cahier des charges rigoureux sur le plan de l'éthique, du professionnalisme, et de la compétence des experts intervenants. Il y trois cadres principaux de certifications sont :

% begin PRZ==========================
\begin{frame}
\frametitle<presentation>{Référentiels ANSSI}
% end header PRZ======================
\begin{itemize}
\item PASSI : Prestataire d’Audit de la sécurité des systèmes d’information ;
\item PDIS : Prestataire de détection d’incident de sécurité ;
\item PRIS : Prestataire de réponse à incident.
\end{itemize}

Ces trois référentiels définissent l’ensemble des exigences d’assurance pour\g {qualifier} des prestataires de services en cybersécurité sur ces trois thématiques. En effet, il serait en effet important de confier la recherche de ses vulnérabilités, leurs remédiations à des sociétés de confiance. 
\end{frame}
% end PRZ==========================
A ces trois volets il ne faut pas oublier, le volet administration des briques informatiques et de télécommunications de l’environnement de l’entreprise. C’est un volet que nous traiterons pas directement dans c\edoc pour se concentrer sur les mécanismes de maintien en continu du niveau de sécurité de l’entreprise avec des mécanismes de veille, d’alerte et de réaction.

%TODO Leure et Hony Posts
%Leurres et Pots de Miel 
%Renseignement

\subsection{Stratégies d'action}

\upicture{Latex/Sources/EDU/Pictures/img-vti}{Les différentes actions de veille}{0.9}{lblvti}

% Begin PRZ ===========================
\begin{frame}
\frametitle<presentation>{Stratégies de l'action}

% end header PRZ =======================
La cyberdéfense est un ensemble de mécanismes liés à une stratégie de l'action. Les outils de cyberdéfense sont construits pour aider à surveiller l'environnement, détecter des menaces et/ou des attaques mais surtout agir et réagir pour limiter les impacts. Si les outils de protection sont configurés à partir d'éléments de politique de sécurité (droits, accès, filtrage ...), les outils de défense sont basés sur les stratégies des attaquants.
On distinguera donc ici trois grands mécanismes de Cyberdefense que les anglo-saxons appellent : 

\begin{itemize}
	\item Predictive Cyberdefense;
	\item Active and Proactive Cyberdefense;
	\item Reactive Cyberdefense.
\end{itemize}
\end{frame}

Nous aborderons, en particulier ces concepts quand nous évoquerons la notion de SOC (Security Operational Center) activité qui opère ce volet de cyberdéfense sachant que la veille sur l'environnement numérique reste un axe important.



% End PRZ ===========================

Il ne faut pas, par ailleurs, oublier le renseignement (\textit{Intelligence}), qui reste une des grandes étapes de la cyberdéfense domaine que nous explorerons sous son volet cyber avec les sources de \g{threat intelligence}, mais aussi avec le Renseignement d'Origine Cyber que les anglo-saxons nomme \g{intelligence cyber}

Dans les grandes organisations, une autre stratégie globale de la cyberdéfense est de penser l'anticipation et la détection de manière globale à l'environnement digital de l'entreprise tout en structurant la réaction de manière locale. 

%TODO https://en.wikipedia.org/wiki/Proactive_cyber_defence

Nous avons positionné l'audit technique comme une des activités fondamentales de la gestion des vulnérabilités.
En effet les techniques d'audit font partie des méthodes de référence pour disposer d'un état des fragilités de l'entreprise. On y trouvera donc les grands basiques des audits techniques que sont les tests d'intrusion, la sécurité applicative, l'audit de configuration, et le fuzzing.

Par ailleurs nous explorerons rapidement, les techniques de déception et de leurre qui font partie cette défense proactive avec les honeypots qui peuvent être couplés avec le \UKword{cyber-hunting}, technique de chasse aux codes malveillants dans l'entreprise.
 

\subsection{Les modèles de cybersécurité}

Les modèles ou framework de sécurité sont intéressants à plusieurs titres. Ils permettent en particulier, de disposer d'un ensemble d'exigences ou de bonnes pratiques organisées dans un référentiel connu et reconnu d'une communauté et  offre un cadre pour : 

\begin{itemize}
  \item Passer en revue les pratiques des organisations en matière de cybersécurité;
  \item Etablir ou améliorer son propre programme de cybersécurité;
  \item Effectuer une auto-évaluation des risques en matière de cybersécurité;
  \item Sensibiliser collaborateurs, partenaires, sous-traitants;
  \item Améliorer la communication entre organisations grâce à un échange d’exigences cybersécurité entre les partenaires commerciaux, les fournisseurs et les régulateurs (ANSSI par exemple)
\end{itemize}


% begin PRZ==========================
\begin{frame}
\frametitle<presentation>{Modèle NIST}
% end header PRZ======================

Il existe de nombreux modèles de description de l'activité de Cyberdefense dans un contexte de cybersécurité.
Certains sont totalement intégrés au modèle de cybersécurité comme l'ISO 27K, ou le Cybersecurity FrameWork du NIST \uref{FrameWork du NIST}{lblnist} avec les activités \textbf{DETECT, RESPOND et RECOVER};
\upicture{Latex/Sources/EDU/Pictures/img-nist}{modèle NIST}{0.9}{lblnist}
\end{frame}
% end PRZ==========================


% begin PRZ==========================
\begin{frame}
\frametitle<presentation>{NIST vs ISO 27k}
% end header PRZ======================
Ce que l'on peut reprocher au modèle du NIST, c'est qu'il ne possède pas explicitement la gestion des fragilités / vulnérabilités, mais il apporte toutefois un modèle très détaillé, que nous utiliserons pour partie.
\upicture{Latex/Sources/EDU/Pictures/img-risk27}{modèle ISO27 et risques}{0.9}{lblrisk27}
\end{frame}
% end PRZ==========================
Dans l'environnement ISO 27000, le modèle est piloté par les risques.\uref{Gestion des risques}{lblrisk27}.

Nous avons fait le choix de positionner la présentation du volet sécurité opérationnelle en nous éloignant un peu des modèles pour présenter les trois grands moteurs de la sécurité opérationnelle. En effet les modèles cités sont orientés sur un axe de cycle de vie.
En sécurité opérationnelle ou cyberdéfense, l'objectif est de conduire en continu et de front des processus de maîtrise des risques cyber opérationnels.
\begin{itemize}
  \item Les systèmes d'information évoluent en continu et des vulnérabilités peuvent s'insérer et/ou être découvertes chaque jour au grès des modifications et évolutions,
  \item Des menaces se concrétisent quotidiennement par des attaques ciblées ou non, nécessitant de réagir vite et en cohérence avec des enjeux de l'entreprise
  \item Avoir la capacité de réagir, et d'assurer la continuité d'activité face à des attaques d'ampleur, ou à fort impact techniques ou médiatique.
\end{itemize}

\newpage
\section{Processus SECOPS}


% Begin PRZ ===========================
\begin{frame}
\frametitle<presentation>{SECOPS en 3 thématiques}
% end header PRZ =======================
Notre propos sera donc centré sur ces trois axes  qui nous déclinerons dans trois chapitres. Le travail de fond d'une équipe de sécurité opérationnelle, ou simplement de l'activité SECOPS est de pouvoir gérer de front trois grandes tâches : 
\begin{itemize}
 \item maîtriser les fragilités numériques de l'entreprise (\UKword{Vulnerability Management)} quelles soient au sein du SI ou dans l'environnement dit digital de cette entreprise (réseaux sociaux, partenaires, ...);
 \item anticiper les menaces et les scénarios associés (\UKword{Threat Management)}, détecter les attaques et gérer au quotidien les événements de sécurité ;
 \item réagir vite et en cohérence avec l'activité de l'entreprise en cas d'incident (\UKword{Incident Management)}.
\end{itemize}
\end{frame}
% end PRZ ===========================

Nous aborderons aussi quelques compléments à ces processus SECOPS, comme la détection des fuites de données (Leak Detection), qui peut s'entendre comme un incident de sécurité externe, ou une détection d'évènements hors de périmètre du système informatique, mais dans le périmètre de surveillance.

% Begin PRZ ===========================
\begin{frame}
\frametitle<presentation>{Les processus SECOPS}
% end header PRZ =======================
\upicture{Latex/Sources/EDU/Pictures/img-chapvar}{Synthèse des méta-processus SECOPS}{0.8}{lblprosecops}

\end{frame}
% end PRZ ===========================

Ces activités nécessitent, pour être efficace, une symbiose parfaite entre les équipes qui gèrent l'activité digitale (Systèmes d'informations, réseaux sociaux,  communication...) et les équipes de sécurité opérationnelle.
Il ne faut pas oublier bien entendu les mécanismes de gouvernance sécurité globale (ISO 27001 par exemple) dans lesquels s'inscrit la sécurité opérationnelle. On trouve souvent dans les entreprises un RSSI dédié cette activité relevant soit du RSSI de la DSI soit d'un DSSI (Directeur de la SSI).

\section{Les métiers de la SECOPS}

Au delà des métiers de l'audit de sécurité qui existent depuis de nombreuses années, la sécurité opérationnelle est le champ de développement de nombreux métiers nouveaux ou en devenir. Nous en explorerons quelques-uns dans chacune des parties qui présentent les opérations de cette SECOPS.

% Begin PRZ ===========================
\begin{frame}
\frametitle<presentation>{Les métiers SECOPS}
% end header PRZ =======================
\upicture{Latex/Sources/EDU/Pictures/img-metiersVTI}{Des métiers SECOPS}{1}{lblmetiersVTI}
\end{frame}
% end PRZ ===========================

\section{Eléments Communs}

Nous trouverez dans chacun des chapitres un petit condensé du processus avec :

\begin{itemize}
  \item Les objectifs;
  \item Les outils;
  \item Les méthodologies et standards;
  \item Les métiers et compétences.
\end{itemize}
}



%$END_UBODY
}%---------------------------------------------------------------------

%************************************************************
% Chargement  du MODELE
%************************************************************

\umainload



